\documentclass[12pt,a4paper]{article}
\usepackage{cmap}
\usepackage[utf8]{inputenc}
\usepackage[T2A]{fontenc}
\usepackage[russian]{babel}
\usepackage{amsmath,amssymb}
\usepackage[left=25mm,right=25mm,top=20mm,bottom=20mm]{geometry}
\usepackage{enumitem}
\usepackage{longtable}
\usepackage{booktabs}

\newcommand{\half}{\tfrac{1}{2}}
\newcommand{\rw}{r_w}
\newcommand{\Rext}{R_{\mathrm{ext}}}
\newcommand{\ppl}{p_0}
\newcommand{\pwf}{p_{\mathrm{wf}}}
\newcommand{\pb}{p_b}
\newcommand{\Tpl}{T_0}
\newcommand{\divr}[1]{\frac{1}{r}\frac{\partial}{\partial r}\!\left(r\,#1\right)}
\newcommand{\pdt}[1]{\frac{\partial #1}{\partial t}}
\newcommand{\pdr}[1]{\frac{\partial #1}{\partial r}}
\newcommand{\csat}{c_o^{\mathrm{sat}}}
\newcommand{\bsat}{\beta_o^{\mathrm{sat}}}
\newcommand{\xio}{\xi_o}
\newcommand{\mdot}{\dot{m}}
\newcommand{\rhocpeff}{(\rho c_p)_{\mathrm{eff}}}
\newcommand{\rhocpconv}{(\rho c_p u)_{\mathrm{conv}}}
\newcommand{\leff}{\lambda_{\mathrm{eff}}}
\newcommand{\QJT}{Q_{\mathrm{JT}}}
\newcommand{\Qad}{Q_{\mathrm{ad}}}
\newcommand{\Qdeg}{Q_{\mathrm{deg}}}

\title{Постановка задачи и численная схема\\
для двухфазной фильтрации Black\,Oil\\
с уравнением энергии}
\author{}
\date{}

\begin{document}
\maketitle
\tableofcontents
\newpage

%=============================================================================
% ВВЕДЕНИЕ (ненумерованное)
%=============================================================================

Настоящий документ содержит полную постановку задачи о двухфазной
фильтрации нефти с растворённым газом (модель Black\,Oil) в призабойной
зоне пласта (ПЗП) с учётом уравнения баланса энергии, а~также описание
численной схемы для решения полученной системы дифференциальных уравнений
в~частных производных.

Документ является самодостаточным: все уравнения, обозначения и~промежуточные
выкладки приведены в~тексте. Конкретные PVT-корреляции (Standing, Beggs--Robinson
и~др.) здесь не~приводятся~--- они реализованы в~программном модуле PVT;
в~документе приводятся лишь требования к~PVT-модели~--- набор функций и~их производных,
которые она должна предоставлять.

\textbf{Структура документа.}
В~разделе~\ref{sec:formulation} описывается физическая постановка: область,
допущения, система уравнений в~непрерывной форме, замыкающие соотношения,
граничные и~начальные условия.
В~разделе~\ref{sec:splitting} формулируется разделение задачи на
гидродинамическую и~тепловую подзадачи и~приводятся упрощённые формулы
коэффициентов.
Раздел~\ref{sec:discretization} посвящён пространственной дискретизации
методом конечных объёмов на логарифмической радиальной сетке.
В~разделе~\ref{sec:time} описывается интегрирование по~времени: метод
Ньютона для совместной системы давления и~насыщенности, алгоритм одного
шага, критерии сходимости.
Приложения содержат подробные выводы промежуточных результатов:
вывод уравнения для темпа выделения свободного газа,
полный вывод уравнения энергии из~первого начала термодинамики,
раскрытие балансов масс и~получение уравнения давления.


%=============================================================================
% ТАБЛИЦА ОБОЗНАЧЕНИЙ
%=============================================================================
\newpage
\section*{Таблица обозначений}
\addcontentsline{toc}{section}{Таблица обозначений}

\begin{longtable}{lll}
\toprule
Символ & Смысл & Ед.~изм. \\
\midrule
\endfirsthead
\toprule
Символ & Смысл & Ед.~изм. \\
\midrule
\endhead
\midrule
\multicolumn{3}{r}{\small\itshape Продолжение на следующей странице} \\
\bottomrule
\endfoot
\bottomrule
\endlastfoot

\multicolumn{3}{l}{\textbf{Координаты и время}} \\
$r$           & радиальная координата                    & м \\
$t$           & время                                     & с \\

\multicolumn{3}{l}{\textbf{Давление и температура}} \\
$p(r,t)$      & давление (единое для обеих фаз)          & Па \\
$T(r,t)$      & температура (единая для породы и флюидов) & К \\
$\pb(T)$      & давление насыщения (bubble point)         & Па \\
$\pwf(t)$     & забойное давление                         & Па \\
$\ppl$        & начальное пластовое давление              & Па \\
$\Tpl$        & начальная пластовая температура            & К \\

\multicolumn{3}{l}{\textbf{Насыщенности}} \\
$S_o(r,t)$    & объёмная доля нефтяной фазы; $S_o+S_g=1$ & --- \\
$S_g(r,t)$    & объёмная доля газовой фазы                & --- \\

\multicolumn{3}{l}{\textbf{Газосодержание}} \\
$R_s^{\text{eq}}(p,T)$ & равновесное газосодержание (максимальная растворимость) & м$^3$ст/м$^3$ст \\
$R_s$         & фактическое содержание газа в нефтяной фазе & м$^3$ст/м$^3$ст \\
$R_s^0$       & начальное газосодержание ($R_s^0 \leqslant R_s^{\text{eq}}(p_0, T_0)$) & м$^3$ст/м$^3$ст \\

\multicolumn{3}{l}{\textbf{PVT-свойства}} \\
$B_o(p,T,R_s)$ & объёмный коэффициент нефти               & м$^3$пл/м$^3$ст \\
$B_g(p,T)$    & объёмный коэффициент газа                 & м$^3$пл/м$^3$ст \\
$\rho_o(p,T,R_s)$ & плотность нефтяной фазы (пласт.~усл.),
                     $\rho_o = (\rho_o^{\text{ст}} + R_s\,\rho_g^{\text{ст}})/B_o$ & кг/м$^3$ \\
$\rho_g(p,T)$ & плотность газовой фазы (пласт.~усл.),
                $\rho_g = \rho_g^{\text{ст}}/B_g$         & кг/м$^3$ \\
$\rho_o^{\text{ст}}$ & плотность нефтяного компонента (ст.~усл., постоянная) & кг/м$^3$ \\
$\rho_g^{\text{ст}}$ & плотность газового компонента (ст.~усл., постоянная)  & кг/м$^3$ \\
$\mu_o(p,T)$  & динамическая вязкость нефти               & Па$\cdot$с \\
$\mu_g(p,T)$  & динамическая вязкость газа                & Па$\cdot$с \\
$c_{p,o}$     & удельная изобарная теплоёмкость нефти     & Дж/(кг$\cdot$К) \\
$c_{p,g}$     & удельная изобарная теплоёмкость газа      & Дж/(кг$\cdot$К) \\
$\alpha_o$    & коэфф.\ объёмного теплового расширения нефти,
                $-(1/\rho_o)\,\partial\rho_o/\partial T|_{p,R_s}$ & К$^{-1}$ \\
$\alpha_g$    & коэфф.\ объёмного теплового расширения газа,
                $-(1/\rho_g)\,\partial\rho_g/\partial T|_p$ & К$^{-1}$ \\
$L(p,T)$      & удельная теплота дегазации                & Дж/кг газа \\

\multicolumn{3}{l}{\textbf{Сжимаемости и коэфф.\ расширения объёмных факторов}} \\
\multicolumn{3}{l}{\quad\small (при фиксированном составе фазы)} \\
$c_o$         & изотермич.\ сжимаемость нефти, $-(1/B_o)\,\partial B_o/\partial p|_{T,R_s}$ & Па$^{-1}$ \\
$c_g$         & изотермич.\ сжимаемость газа, $-(1/B_g)\,\partial B_g/\partial p|_T$       & Па$^{-1}$ \\
$\beta_o$     & коэфф.\ теплового расширения $B_o$, $(1/B_o)\,\partial B_o/\partial T|_{p,R_s}$ & К$^{-1}$ \\
$\beta_g$     & коэфф.\ теплового расширения $B_g$, $(1/B_g)\,\partial B_g/\partial T|_p$     & К$^{-1}$ \\
$\csat$       & полная сжимаемость нефти вдоль кривой насыщения & Па$^{-1}$ \\
$\bsat$       & полный коэфф.\ расширения $B_o$ вдоль кривой насыщения & К$^{-1}$ \\

\multicolumn{3}{l}{\textbf{Пористая среда}} \\
$\phi$        & пористость (постоянная)                   & --- \\
$k$           & абсолютная проницаемость (постоянная)     & м$^2$ \\
$c_r$         & сжимаемость порового пространства ($=0$ в данной модели) & Па$^{-1}$ \\
$k_{ro}(S_o)$ & относит.\ фазовая проницаемость нефти    & --- \\
$k_{rg}(S_g)$ & относит.\ фазовая проницаемость газа     & --- \\
$(\rho c_p)_s$& объёмная теплоёмкость скелета породы (индекс $s$~--- solid) & Дж/(м$^3\cdot$К) \\
$\leff$       & теплопроводность гомогенной среды (порода\,+\,нефть\,+\,газ) & Вт/(м$\cdot$К) \\
$\hat\lambda_s$   & теплопроводность скелета породы       & Вт/(м$\cdot$К) \\
$\hat\lambda_o$   & теплопроводность нефтяной фазы         & Вт/(м$\cdot$К) \\
$\hat\lambda_g$   & теплопроводность газовой фазы           & Вт/(м$\cdot$К) \\

\multicolumn{3}{l}{\textbf{Термодинамика}} \\
$h_o$, $h_g$  & удельные энтальпии фаз                    & Дж/кг \\
$e_o$, $e_g$  & удельные внутренние энергии фаз           & Дж/кг \\
$\bar{h}_{oo}$& парциальная энтальпия нефтяного комп.\ в~нефтяной фазе & Дж/кг \\
$\bar{h}_{go}$& парциальная энтальпия газового комп.\ в~нефтяной фазе  & Дж/кг \\
$\mdot$       & массовый темп выделения свободного газа   & кг/(м$^3\cdot$с) \\

\multicolumn{3}{l}{\textbf{Скорости и подвижности}} \\
$u_o$, $u_g$  & скорости фильтрации (Дарси)              & м/с \\
$\Lambda_\alpha$ & подвижность фазы $\alpha$, $k\,k_{r\alpha}/\mu_\alpha$ & м$^2$/(Па$\cdot$с) \\

\multicolumn{3}{l}{\textbf{Вспомогательные}} \\
$\xio$        & $S_o/B_o = (1-S_g)/B_o$                  & м$^3$ст/м$^3$пор \\
$\xi_g$       & $R_s\,\xio + S_g/B_g$                    & м$^3$ст/м$^3$пор \\

\multicolumn{3}{l}{\textbf{Геометрия}} \\
$\rw$         & радиус скважины                           & м \\
$\Rext$       & радиус контура питания                    & м \\
$h$           & мощность (толщина) пласта                  & м \\

\end{longtable}


%=============================================================================
\section{Физическая постановка}\label{sec:formulation}
%=============================================================================

%-----------------------------------------------------------------------------
\subsection{Область и допущения}\label{ssec:domain}
%-----------------------------------------------------------------------------

Рассматривается горизонтальный пласт постоянной мощности~$h$, вскрытый
вертикальной скважиной радиуса~$\rw$. Пласт ограничен контуром
питания~$r = \Rext$, через который поступает жидкость из удалённой зоны.
Задача одномерная~--- все величины зависят только от радиальной
координаты~$r$ и~времени~$t$:
\begin{equation}\label{eq:domain}
  r \in [\rw,\; \Rext], \qquad t \geqslant 0.
\end{equation}

Оператор дивергенции в цилиндрических координатах при осевой
симметрии имеет вид
\begin{equation}\label{eq:div}
  \mathrm{div}\,\mathbf{q} = \divr{q_r},
\end{equation}
где~$q_r$~--- радиальная компонента потока.

Пластовый флюид описывается в~рамках модели Black\,Oil с~двумя фазами:
\begin{itemize}[nosep]
  \item \textbf{нефтяная фаза}~(индекс~$o$)~--- жидкая смесь тяжёлого
        нефтяного компонента и растворённого в~нём лёгкого газового
        компонента;
  \item \textbf{газовая фаза}~(индекс~$g$)~--- свободный газ, выделившийся
        из нефти при снижении давления ниже давления насыщения.
\end{itemize}

В~модели Black\,Oil различают \emph{компоненты} (нефтяной и~газовый~---
химические вещества) и~\emph{фазы} (нефтяная и~газовая~--- агрегатные
состояния). Нефтяной компонент содержится только в~нефтяной фазе.
Газовый компонент может присутствовать в~обеих фазах: в~нефтяной~---
в~растворённом виде, в~газовой~--- в~свободном. Водная фаза отсутствует.
Испарение нефтяного компонента в~газовую фазу не~учитывается: газовая
фаза содержит только газовый компонент.

\paragraph{Допущения.}

\begin{enumerate}[label=\arabic*., nosep]
  \item \textbf{Мгновенное локальное фазовое равновесие.}
        При $p < \pb(T)$ газ выделяется из нефти мгновенно,
        чтобы текущее газосодержание опустилось до равновесного значения.
        Кинетика фазового перехода не учитывается.

        Различаются три величины:
        \begin{itemize}[nosep]
          \item $R_s^{\text{eq}}(p,T)$~--- \textbf{равновесное
                газосодержание}: максимальное количество газа, которое
                может быть растворено в~нефти при данных~$p$ и~$T$.
                Функция, предоставляемая PVT-моделью.
          \item $R_s(r,t)$~--- \textbf{фактическое газосодержание}
                нефтяной фазы. Является \textbf{неизвестной} задачи.
                Всегда $R_s \leqslant R_s^{\text{eq}}(p,T)$.
          \item $R_s^0$~--- \textbf{начальное газосодержание} нефти
                в~пласте; $R_s^0 \leqslant R_s^{\text{eq}}(\ppl, \Tpl)$;
                равенство соответствует нефти, насыщенной газом при
                начальных условиях.
        \end{itemize}

        В~двухфазной зоне, $S_g > 0$, нефть насыщена, и~фазовое
        равновесие требует $R_s = R_s^{\text{eq}}(p,T)$. В~однофазной
        зоне, $S_g = 0$, фактическое газосодержание определяется из
        баланса газового компонента и~зависит от~истории процесса.
        При простых режимах эксплуатации (монотонное снижение давления,
        отсутствие дополнительных источников газа, изотермическая
        гидродинамика) возможны ситуации, когда $R_s = R_s^0 = \mathrm{const}$
        в~однофазной зоне, однако в~общем случае конкретной формулы
        для~$R_s$ нет, и оно определяется решением задачи.

  \item \textbf{Единое давление фаз.}
        Капиллярное давление не~учитывается: $p_o = p_g \equiv p$.
        Допустимо при высокой проницаемости и~умеренных
        газонасыщенностях.
  \item \textbf{Температурная зависимость PVT-свойств.}
        Все свойства фаз~--- объёмные коэффициенты $B_o$, $B_g$,
        плотности $\rho_o$, $\rho_g$, вязкости $\mu_o$, $\mu_g$,
        коэффициенты теплового расширения $\alpha_o$, $\alpha_g$,
        теплоёмкости $c_{p,o}$, $c_{p,g}$~--- зависят как от~давления,
        так и~от~температуры.
        В~гидродинамической подзадаче (раздел~\ref{sec:splitting}) эта
        зависимость упрощается.
  \item \textbf{Локальное термодинамическое равновесие порода--флюид.}
        Температура породы и~флюидов одинакова в~каждой точке:
        $T_s = T_f = T$.
  \item \textbf{Постоянная пористость.}
        $\phi = \mathrm{const}$, сжимаемость породы $c_r = 0$.
        В~разделе~\ref{ssec:equations} формулы записываются в~общем виде
        при $c_r \neq 0$ для полноты; при упрощении
        (раздел~\ref{sec:splitting}) $c_r = 0$.
  \item \textbf{Отсутствие теплообмена с~кровлей и~подошвой.}
        Задача осреднена по~вертикали; тепловые потери через кровлю
        и~подошву пласта не~учитываются. Все величины ($p$, $T$, $S_g$,
        $R_s$) представляют собой средние по~толщине пласта~$h$.
\end{enumerate}

\paragraph{Неизвестные.}
Задача содержит \textbf{четыре неизвестных}: давление~$p(r,t)$,
газонасыщенность~$S_g(r,t)$, фактическое газосодержание~$R_s(r,t)$
и~температуру~$T(r,t)$.



%-----------------------------------------------------------------------------
\subsection{Система уравнений}\label{ssec:equations}
%-----------------------------------------------------------------------------

Задача содержит \textbf{четыре неизвестных}: давление~$p(r,t)$,
газонасыщенность~$S_g(r,t)$, фактическое газосодержание~$R_s(r,t)$
и~температуру~$T(r,t)$.

%............................................................................
\subsubsection{Закон Дарси}

Скорости фильтрации для каждой фазы:
\begin{equation}\label{eq:darcy}
  \boxed{
    u_\alpha = -\frac{k\,k_{r\alpha}}{\mu_\alpha}\,\pdr{p},
    \qquad \alpha \in \{o, g\}.
  }
\end{equation}
Подвижность фазы:
$\Lambda_\alpha = k\,k_{r\alpha}/\mu_\alpha$,
\quad $u_\alpha = -\Lambda_\alpha\,\partial p/\partial r$.


%............................................................................
\subsubsection{Баланс массы нефтяного компонента}\label{sssec:mass_oil}

\begin{equation}\label{eq:mass_oil}
  \boxed{
    \pdt{}\!\left(\phi\,\frac{1 - S_g}{B_o}\right)
    + \divr{\frac{u_o}{B_o}} = 0.
  }
\end{equation}


%............................................................................
\subsubsection{Баланс массы газового компонента}\label{sssec:mass_gas}

\begin{equation}\label{eq:mass_gas}
  \boxed{
    \pdt{}\!\left(\phi\!\left[
      \frac{R_s\,(1 - S_g)}{B_o} + \frac{S_g}{B_g}
    \right]\right)
    + \divr{\left[
      \frac{R_s\,u_o}{B_o} + \frac{u_g}{B_g}
    \right]} = 0.
  }
\end{equation}

Уравнения~\eqref{eq:mass_oil} и~\eqref{eq:mass_gas} записаны
в~\textbf{дивергентной (консервативной)} форме; именно в~этой форме
они дискретизируются. Раскрытие производных приведено
в~приложении~\ref{app:expansion}.


%............................................................................
\subsubsection{Замыкающее уравнение (переключение режимов)}
\label{sssec:variable_switching}

Третье уравнение связывает~$S_g$ и~$R_s$; его вид зависит
от~\textbf{режима} ячейки.

\paragraph{Режим~1: двухфазная зона} ($S_g > 0$).
Нефть насыщена газом, фазовое равновесие:
\begin{equation}\label{eq:regime1}
  \boxed{R_s = R_s^{\text{eq}}(p).}
\end{equation}

\paragraph{Режим~2: однофазная зона} ($S_g = 0$).
Газовая фаза отсутствует:
\begin{equation}\label{eq:regime2}
  \boxed{S_g = 0.}
\end{equation}
Фактическое~$R_s$ определяется из баланса газового
компонента~\eqref{eq:mass_gas}.


%............................................................................
\subsubsection{Уравнение энергии}\label{sssec:energy}

\begin{equation}\label{eq:main_energy_final}
  \boxed{
    \rhocpeff\,\pdt{T}
    + \rhocpconv\,\pdr{T}
    = \divr{\leff\,\pdr{T}}
    + \QJT + \Qad + \Qdeg.
  }
\end{equation}

Коэффициенты при производных температуры:
\begin{equation}\label{eq:main_rhocpeff}
\rhocpeff
    = (1-\phi)\,(\rho c_p)_s
      + \phi\,(S_o\,\rho_o\,c_{p,o} + S_g\,\rho_g\,c_{p,g}),
\end{equation}
\begin{equation}\label{eq:main_rhocpconv}
\rhocpconv
    = \rho_o\,c_{p,o}\,u_o + \rho_g\,c_{p,g}\,u_g.
\end{equation}

Источниковые члены:
\begin{equation}\label{eq:main_Qad}
\Qad = \phi\left[T\,\sum_\alpha S_\alpha\alpha_\alpha
       + (\rho c_p)_s\,T\,c_r + p\,c_r\right]\pdt{p},
\end{equation}
\begin{equation}\label{eq:main_QJT}
\QJT = -\left[u_g\,(1 - \alpha_g\,T)
       + u_o\,(1 - \alpha_o\,T)\right]\pdr{p},
\end{equation}
\begin{equation}\label{eq:main_Qdeg}
\Qdeg = - L(p,T)\,\mdot,
\end{equation}
где~$L(p,T)$~--- удельная теплота дегазации [Дж/кг газа],
параметр PVT-модели (раздел~\ref{ssec:closures}).

Массовый темп выделения свободного газа:
\begin{equation}\label{eq:main_mdot}
  \boxed{
    \mdot = -\rho_g^{\text{ст}}\!\left[
      \phi\,\frac{1-S_g}{B_o}\,\pdt{R_s}
      + \frac{u_o}{B_o}\,\pdr{R_s}
    \right].
  }
\end{equation}
При~$\mdot > 0$ газ выделяется; при~$\mdot < 0$~--- растворяется.

Тепловую задачу можно записать через коэффициент Джоуля--Томсона
$\mu_{\mathrm{JT},\alpha}$, который заменяет~$\alpha_\alpha$:
\begin{equation}\label{eq:main_muJT_def}
  \mu_{\mathrm{JT},\alpha}
  = \frac{v_\alpha\,(\alpha_\alpha\,T - 1)}{c_{p,\alpha}},
  \qquad
  1 - \alpha_\alpha\,T
  = -\rho_\alpha\,c_{p,\alpha}\,\mu_{\mathrm{JT},\alpha}.
\end{equation}
Тогда
\begin{equation}\label{eq:main_QJT_muJT}
  \boxed{
    \QJT
    = \sum_\alpha \rho_\alpha\,c_{p,\alpha}\,
      \mu_{\mathrm{JT},\alpha}\,u_\alpha\,\pdr{p},
  }
\end{equation}
\begin{equation}\label{eq:main_Qad_muJT}
  \boxed{
    \Qad
    = \phi\!\left[ 1+
      \sum_\alpha S_\alpha\,
        \rho_\alpha\,c_{p,\alpha}\,\mu_{\mathrm{JT},\alpha}
      + c_r\!\left(p + (\rho c_p)_s\,T\right)
    \right]\pdt{p}.
  }
\end{equation}

%............................................................................
\subsubsection{Сводная таблица системы уравнений}

Полная замкнутая система состоит из \textbf{четырёх уравнений}
на~\textbf{четыре неизвестных}: $p(r,t)$, $S_g(r,t)$, $R_s(r,t)$,
$T(r,t)$. Балансы масс записаны в~дивергентной (консервативной)
форме и~дискретизируются непосредственно.

\paragraph{1. Баланс массы нефтяного компонента~\eqref{eq:mass_oil}:}
\[
  \pdt{}\!\left(\phi\,\frac{S_o}{B_o}\right)
  + \divr{\frac{u_o}{B_o}} = 0.
\]

\paragraph{2. Баланс массы газового компонента~\eqref{eq:mass_gas}:}
\[
  \pdt{}\!\left(\phi\!\left[
    \frac{R_s\,S_o}{B_o} + \frac{S_g}{B_g}
  \right]\right)
  + \divr{\left[
    \frac{R_s\,u_o}{B_o} + \frac{u_g}{B_g}
  \right]} = 0.
\]

\paragraph{3. Замыкающее уравнение (переключаемое,
раздел~\ref{sssec:variable_switching}):}
\[
  \text{Режим~1 ($S_g > 0$):}\quad R_s = R_s^{\text{eq}}(p,T);
  \qquad
  \text{Режим~2 ($S_g = 0$):}\quad S_g = 0.
\]

\paragraph{4. Уравнение энергии~\eqref{eq:main_energy_final}:}
\[
  \rhocpeff\,\pdt{T}
  + \rhocpconv\,\pdr{T}
  = \divr{\leff\,\pdr{T}}
  + \QJT + \Qad + \Qdeg.
\]

\paragraph{Неизвестные и~порядок решения:}
$p$, $S_g$, $R_s$~--- определяются из уравнений~1--3
(совместная система, метод Ньютона);
$T$~--- из уравнения~4 на~известных полях $p$, $S_g$, $R_s$.

\paragraph{Замыкающие соотношения:}
\begin{itemize}[nosep]
  \item $S_o = 1 - S_g$ (алгебраическое);
  \item закон Дарси~\eqref{eq:darcy}:
        $u_\alpha = -\Lambda_\alpha\,\partial p/\partial r$;
  \item PVT-корреляции: $R_s^{\text{eq}}(p,T)$, $B_o(p,T,R_s)$,
        $B_g(p,T)$, $\rho_\alpha(p,T)$, $\mu_\alpha(p,T)$,
        $c_{p,\alpha}(p,T)$, $\alpha_\alpha(p,T)$;
  \item теплота дегазации $L(p,T)$;
  \item ОФП: $k_{ro}(S_o)$, $k_{rg}(S_g)$;
  \item свойства породы: $\phi$, $k$, $(\rho c_p)_s$;
  \item теплопроводность среды~\eqref{eq:lambda_eff}:
        $\leff = (1{-}\phi)\,\hat\lambda_s
        + \phi\,(S_o\,\hat\lambda_o + S_g\,\hat\lambda_g)$.
\end{itemize}

\paragraph{Граничные условия:}
\begin{itemize}[nosep]
  \item $p(\rw,t) = \pwf(t)$, \quad $p(\Rext,t) = \ppl$;
  \item $T(\Rext,t) = \Tpl$, \quad
        $\partial T/\partial r\big|_{\rw} = 0$;
  \item $S_g(\Rext,t) = 0$, \quad $R_s(\Rext,t) = R_s^0$.
\end{itemize}

\paragraph{Начальные условия:}
$p(r,0) = \ppl$, \quad $T(r,0) = \Tpl$, \quad
$S_g(r,0) = 0$, \quad $R_s(r,0) = R_s^0$.


%=============================================================================
\subsection{Замыкающие соотношения}\label{ssec:closures}
%=============================================================================

\subsubsection{Требования к PVT-модели}

Конкретные корреляции (Standing, Beggs--Robinson, Lee--Gonzalez--Eakin
и~др.) здесь не~приводятся~--- они реализованы в~программном модуле PVT.
Программная реализация предусматривает возможность замены корреляций
без изменения солвера. Ниже приведены требования~--- набор функций
и~их производных, которые должна предоставлять любая PVT-модель.

\begin{center}
\begin{tabular}{llll}
\toprule
Функция & Аргументы & Ед.~изм. & Производные для солвера \\
\midrule
$R_s^{\text{eq}}$ & $p, T$ & м$^3$ст/м$^3$ст
  & $\partial R_s^{\text{eq}}/\partial p$, $\partial R_s^{\text{eq}}/\partial T$ \\
$B_o$     & $p, T, R_s$ (фактическое) & м$^3$пл/м$^3$ст
  & $\partial B_o/\partial p|_{T,R_s}$,
    $\partial B_o/\partial T|_{p,R_s}$,
    $\partial B_o/\partial R_s|_{p,T}$ \\
$B_g$     & $p, T$ & м$^3$пл/м$^3$ст
  & $\partial B_g/\partial p|_T$ \\
$\rho_o$  & $p, T$ & кг/м$^3$
  & через $B_o$, $R_s$~\eqref{eq:rho_oil} \\
$\rho_g$  & $p, T$ & кг/м$^3$
  & $\partial\rho_g/\partial p$ \\
$\mu_o$   & $p, T$ & Па$\cdot$с
  & $\partial\mu_o/\partial p$ \\
$\mu_g$   & $p, T$ & Па$\cdot$с
  & $\partial\mu_g/\partial p$ \\
$c_{p,o}$, $c_{p,g}$ & $p, T$ & Дж/(кг$\cdot$К)
  & --- \\
$\alpha_o$, $\alpha_g$ & $p, T$ & К$^{-1}$
  & --- \\
$L$       & $p, T$ & Дж/кг газа
  & --- \\
$\pb$     & $T$ & Па
  & --- \\
\bottomrule
\end{tabular}
\end{center}


\subsubsection{Сжимаемости и коэффициенты расширения}

При \textbf{фиксированном составе} (для уравнения энергии):
\begin{equation}\label{eq:co_cg}
  c_o = -\frac{1}{B_o}\left.\frac{\partial B_o}{\partial p}\right|_{T,R_s},
  \qquad
  c_g = -\frac{1}{B_g}\left.\frac{\partial B_g}{\partial p}\right|_T,
\end{equation}
\begin{equation}\label{eq:beta_o_g}
  \beta_o = \frac{1}{B_o}\left.\frac{\partial B_o}{\partial T}\right|_{p,R_s},
  \qquad
  \beta_g = \frac{1}{B_g}\left.\frac{\partial B_g}{\partial T}\right|_p.
\end{equation}

Условие $R_s = \mathrm{const}$ означает фиксированное
\textbf{фактическое} содержание газа в~нефтяной фазе~--- это
принципиально для нефтяной фазы.
Полная производная по~$T$ при фиксированном~$p$ включает вклад
изменения состава:
\begin{equation}\label{eq:beta_o_full}
  \left.\frac{\partial B_o}{\partial T}\right|_p
  = \left.\frac{\partial B_o}{\partial T}\right|_{p,R_s}
    + \left.\frac{\partial B_o}{\partial R_s}\right|_{p,T}
      \left.\frac{\partial R_s}{\partial T}\right|_p.
\end{equation}
В~уравнении энергии используется первое слагаемое ($\beta_o$);
вклад изменения состава учтён через~$\Qdeg$.

\textbf{Вдоль кривой насыщения} (для балансов масс)~---
$\csat$~\eqref{eq:co_sat} и~$\bsat$~\eqref{eq:beta_o_sat}.

Коэффициент теплового расширения фазы (при фиксированном составе):
\begin{equation}\label{eq:thermal_expansion}
  \alpha_\alpha = -\frac{1}{\rho_\alpha}
    \left.\frac{\partial\rho_\alpha}{\partial T}\right|_{p,R_s},
    \qquad \alpha \in \{o, g\}.
\end{equation}


\subsubsection{Требования к модели ОФП}

Модель ОФП предоставляет функции $k_{ro}(S_g)$ и~$k_{rg}(S_g)$
и~их производные~$dk_{ro}/dS_g$, $dk_{rg}/dS_g$ (необходимы для
блочного якобиана совместной системы).

Модель Кори~--- частный случай реализации. Интерфейс допускает
табличные ОФП, модель Брукса--Кори, LET и~др.

Пример~--- \textbf{модель Кори}:
нормализованная насыщенность $S^* = (S_g - S_{gc})/(1 - S_{or} - S_{gc})$;
\begin{equation}\label{eq:corey}
  k_{ro} = k_{ro,\max}\,(1-S^*)^{n_o}, \qquad
  k_{rg} = k_{rg,\max}\,(S^*)^{n_g}.
\end{equation}
Параметры: $S_{or}$, $S_{gc}$, $n_o$, $n_g$, $k_{ro,\max}$, $k_{rg,\max}$.


\subsubsection{Свойства породы}

\begin{itemize}[nosep]
  \item $\phi$~--- пористость, постоянная.
  \item $k$~--- абсолютная проницаемость (постоянная), м$^2$.
  \item $(\rho c_p)_s$~--- объёмная теплоёмкость скелета,
        Дж/(м$^3\cdot$К).
\end{itemize}


\subsubsection{Эффективная теплопроводность}

Параллельная (объёмно-взвешенная) модель:
\begin{equation}\label{eq:lambda_eff}
  \boxed{
    \leff = (1-\phi)\,\hat\lambda_s + \phi\,(S_o\,\hat\lambda_o + S_g\,\hat\lambda_g).
  }
\end{equation}
Параметры $\hat\lambda_s$, $\hat\lambda_o$, $\hat\lambda_g$~--- постоянные,
задаются в~конфигурационном файле. $\leff$ изменяется по~пласту вместе
с~$S_g$.


%=============================================================================
\subsection{Граничные и начальные условия}\label{ssec:bc_ic}
%=============================================================================

\subsubsection{Начальные условия}\label{sssec:ic}

В~момент $t = 0$ пласт находится в~однородном состоянии:
\begin{equation}\label{eq:ic}
  p(r, 0) = \ppl, \qquad T(r, 0) = \Tpl, \qquad S_g(r, 0) = 0.
\end{equation}
Начальное давление $\ppl > \pb(\Tpl)$: весь газ растворён в~нефти,
свободная газовая фаза отсутствует.


\subsubsection{Граничные условия по давлению}\label{sssec:bc_p}

\paragraph{На забое скважины ($r = \rw$):}
\begin{equation}\label{eq:bc_p_well}
  p(\rw, t) = \pwf(t),
\end{equation}
где $\pwf(t)$~--- заданная функция времени (забойное давление).
Забойное давление задаётся \textbf{историей}: набор пар
(интервал времени, значение давления). На~каждом интервале
$\pwf = \mathrm{const}$, между интервалами~--- скачок:
\begin{equation}\label{eq:pwf_history}
  \pwf(t) = \pwf^{(k)} \quad \text{при}\quad t \in [t_k,\; t_{k+1}).
\end{equation}
При типичных условиях $\pwf < \pb$, поэтому вблизи скважины
образуется двухфазная зона.

\paragraph{На контуре питания ($r = \Rext$):}
\begin{equation}\label{eq:bc_p_ext}
  p(\Rext, t) = \ppl = \mathrm{const}.
\end{equation}
Через контур питания жидкость поступает в~пласт при постоянном давлении.


\subsubsection{Граничные условия по температуре}\label{sssec:bc_T}

\paragraph{На контуре питания ($r = \Rext$):}
\begin{equation}\label{eq:bc_T_ext}
  T(\Rext, t) = \Tpl.
\end{equation}
Приток с~контура имеет пластовую температуру.

\paragraph{На забое скважины ($r = \rw$):}
\begin{equation}\label{eq:bc_T_well}
  \left.\pdr{T}\right|_{r = \rw} = 0
\end{equation}
--- условие Неймана (нулевой кондуктивный поток). Температура
у~скважины определяется конвекцией, а~не~теплопроводностью.
Кондуктивный поток в~скважину отсутствует.


\subsubsection{Граничные условия по насыщенности и~газосодержанию}
\label{sssec:bc_Sg}

\paragraph{На контуре питания ($r = \Rext$):}
\begin{equation}\label{eq:bc_Sg_ext}
  S_g(\Rext, t) = 0, \qquad R_s(\Rext, t) = R_s^0.
\end{equation}
Нефть поступает через контур питания в~однофазном состоянии
($p = \ppl > \pb$, свободного газа нет) с~постоянным
газосодержанием, равным начальному.

\paragraph{На забое ($r = \rw$):}
граничные условия для~$S_g$ и~$R_s$ \textbf{не задаются}. Давление
на~забое фиксировано ($p = \pwf$). Газонасыщенность и~газосодержание
на~забое~--- свободные переменные, определяемые динамикой дегазации
и~балансами масс.


%=============================================================================
\section{Разделение задачи}\label{sec:splitting}
%=============================================================================

%-----------------------------------------------------------------------------
\subsection{Допущение независимости гидродинамики от температуры}
\label{ssec:T_independent}
%-----------------------------------------------------------------------------

Все PVT-свойства в~уравнениях баланса масс~\eqref{eq:mass_oil},
\eqref{eq:mass_gas} считаются зависящими
\textbf{только от~давления} при фиксированной температуре
$T = \Tpl = \mathrm{const}$:
\[
  R_s(p, \Tpl), \quad B_o(p, \Tpl), \quad B_g(p, \Tpl), \quad
  \mu_o(p, \Tpl), \quad \mu_g(p, \Tpl).
\]

\paragraph{Следствия:}
\begin{itemize}[nosep]
  \item $\partial R_s/\partial T = 0$, $\partial B_o/\partial T = 0$,
        $\partial B_g/\partial T = 0$ в~гидродинамических уравнениях;
  \item $\beta_o = 0$, $\beta_g = 0$, $\bsat = 0$
        в~гидродинамических уравнениях;
  \item в~уравнениях баланса масс все зависимости от~$T$ исключены;
        гидродинамическая подсистема определяет $p$, $S_g$ и~$R_s$.
\end{itemize}


В~уравнении энергии PVT-свойства сохраняют полную
зависимость от~$p$ и~$T$: $c_{p,\alpha}(p,T)$, $\alpha_\alpha(p,T)$,
$L(p,T)$ и~т.\,д.


%-----------------------------------------------------------------------------
\subsection{Упрощение гидродинамических уравнений}\label{ssec:simplification}
%-----------------------------------------------------------------------------

При $T = \Tpl = \mathrm{const}$ и~$c_r = 0$ в~уравнениях
баланса масс~\eqref{eq:mass_oil}, \eqref{eq:mass_gas} упрощаются
PVT-свойства:
\begin{itemize}[nosep]
  \item $R_s$, $B_o$, $B_g$, $\mu_o$, $\mu_g$ зависят только от~$p$;
  \item все производные по~$T$ обращаются в~нуль;
  \item $\phi = \mathrm{const}$ и~выносится из-под производной по~$t$.
\end{itemize}

Консервативные уравнения сохраняют свою форму:
\begin{equation}\label{eq:mass_oil_simplified}
  \boxed{
    \phi\,\pdt{}\!\left(\frac{1 - S_g}{B_o(p)}\right)
    + \divr{\frac{u_o}{B_o(p)}} = 0,
  }
\end{equation}
\begin{equation}\label{eq:mass_gas_simplified}
  \boxed{
    \phi\,\pdt{}\!\left(
      \frac{R_s(p)\,(1 - S_g)}{B_o(p)} + \frac{S_g}{B_g(p)}
    \right)
    + \divr{\left[
      \frac{R_s(p)\,u_o}{B_o(p)} + \frac{u_g}{B_g(p)}
    \right]} = 0.
  }
\end{equation}

Неизвестные гидродинамической подсистемы~--- $p(r,t)$, $S_g(r,t)$
и~$R_s(r,t)$; замыкающее уравнение
(раздел~\ref{sssec:variable_switching}) дополняет систему.
Все PVT-свойства~--- функции только~$p$;
ОФП~--- функции только~$S_g$. Уравнения дискретизируются
в~консервативной форме непосредственно
(раздел~\ref{sec:discretization}).

\subsubsection{Уравнение энергии на известных полях}

На~шаге теплопереноса поля $p^{n+1}(r)$, $S_g^{n+1}(r)$, $R_s^{n+1}(r)$
и~скорости $u_o^{n+1}(r)$, $u_g^{n+1}(r)$ уже определены
из~гидродинамической подсистемы. Коэффициенты уравнения
энергии~\eqref{eq:main_energy_final}, зависящие от~$T$
($\rhocpeff$, $\rhocpconv$, $\leff$), замораживаются на~$T^n$.
После замораживания уравнение \textbf{линейно} по~$T^{n+1}$
и~приводится к~виду
\begin{equation}\label{eq:energy_linear}
  A(r)\,\pdt{T} + B(r)\,\pdr{T}
  = \divr{\leff\,\pdr{T}} + f(r),
\end{equation}
где
\begin{align}
  A(r) &= \rhocpeff\big|_{T^n},
    \label{eq:coeff_A} \\
  B(r) &= \rhocpconv\big|_{T^n}
    = \rho_o\,c_{p,o}\,u_o + \rho_g\,c_{p,g}\,u_g,
    \label{eq:coeff_B} \\
  f(r) &= \QJT + \Qad + \Qdeg.
    \label{eq:coeff_f}
\end{align}
Все величины в~$A$, $B$, $f$ вычисляются при~$p^{n+1}$,
$S_g^{n+1}$, $u^{n+1}$ и~$T^n$.

Массовый темп $\mdot$~\eqref{eq:main_mdot} содержит только
производные $R_s(r,t)$ по~пространству и~времени и~не~зависит
от~$T$. Поэтому $\Qdeg = -L\,\mdot$~--- \textbf{известная}
скалярная величина в~каждой ячейке, полностью входящая
в~правую часть~$f(r)$.

После пространственной дискретизации
уравнение~\eqref{eq:energy_linear} даёт \textbf{трёхдиагональную
СЛАУ} относительно~$T_j^{n+1}$, решаемую прогонкой
(алгоритм Томаса) без итераций.


%-----------------------------------------------------------------------------
\subsection{Порядок решения}\label{ssec:order}
%-----------------------------------------------------------------------------

На~каждом временном шаге $t^n \to t^{n+1}$:

\begin{enumerate}
  \item \textbf{Гидродинамика (изотермическая).}
    Решить совместную систему балансов
    масс~\eqref{eq:mass_oil_simplified},
    \eqref{eq:mass_gas_simplified} с~замыкающим
    уравнением (раздел~\ref{sssec:variable_switching})
    методом Ньютона.
    PVT-свойства вычисляются при~$T = \Tpl$ и~обновляются
    на~каждой итерации при текущем приближении~$p^{(k)}$.
    Результат: $p^{n+1}$, $S_g^{n+1}$, $R_s^{n+1}$.
    Скорости Дарси $u_o^{n+1}$, $u_g^{n+1}$ вычисляются
    из~$p^{n+1}$ по~\eqref{eq:darcy}.

  \item \textbf{Теплоперенос.}
    Вычислить коэффициенты $A$, $B$, $\leff$, $f$
    по~\eqref{eq:coeff_A}--\eqref{eq:coeff_f},
    используя $p^{n+1}$, $u^{n+1}$, $S_g^{n+1}$, $R_s^{n+1}$
    и~$T^n$.
    Решить линейную систему~\eqref{eq:energy_linear}
    $\to$ $T^{n+1}$.
\end{enumerate}

Обратная связь $T \to$ гидродинамика \textbf{отсутствует}. Температура~---
пассивный скаляр, переносимый на~известных полях давления
и~насыщенности. Нелинейность уравнения энергии проявляется только
через зависимость коэффициентов от~$T$ (слабая), которая
обрабатывается замораживанием на~$T^n$.


%=============================================================================
\section{Пространственная дискретизация}\label{sec:discretization}
%=============================================================================

%-----------------------------------------------------------------------------
\subsection{Радиальная сетка}\label{ssec:grid}
%-----------------------------------------------------------------------------

Используется метод конечных объёмов (FVM) на логарифмической сетке по~$r$.

\textbf{Сопряжённая сетка}: грани контрольных объёмов $r_{j+\half}$
(сопряжённые узлы) и~центры контрольных объёмов~$r_j$ (основные узлы).

\paragraph{Логарифмическое распределение граней:}
\begin{equation}\label{eq:log_grid}
  \ln\frac{r_{j+\half}}{r_{j-\half}} = \mathrm{const},
  \qquad
  r_{j+\half} = \rw \cdot \left(\frac{\Rext}{\rw}\right)^{j/N},
  \quad j = 0, 1, \ldots, N.
\end{equation}
Крайние грани совпадают с~границами области:
$r_{\half} = \rw$, $r_{N+\half} = \Rext$.

\paragraph{Центры ячеек}
--- центр масс кольцевого слоя (логарифмически взвешенное
среднее граней). Для всех ячеек ($j = 0, 1, \ldots, N-1$):
\begin{equation}\label{eq:cell_center}
  r_j = r_{j+\half}\,
  \sqrt{\exp\!\left(
    -\frac{x\ln x}{1 - x} - 1
  \right)},
  \qquad
  x = \left(\frac{r_{j-\half}}{r_{j+\half}}\right)^2.
\end{equation}
Центры граничных ячеек ($j = 0$ и~$j = N{-}1$) определяются
по~той же формуле, что и~внутренние. Граничные условия
реализуются через потоки на~граничных гранях
(раздел~\ref{ssec:bc_disc}).

\paragraph{Объём ячейки} (с~учётом мощности пласта~$h$):
\begin{equation}\label{eq:volume}
  V_j = \pi\,h\,(r_{j+\half}^2 - r_{j-\half}^2).
\end{equation}

Параметры сетки ($\rw$, $\Rext$, $N$) задаются пользователем
в~конфигурационном файле.


%-----------------------------------------------------------------------------
\subsection{Балансы масс: метод Ньютона}\label{ssec:pressure_disc}
%-----------------------------------------------------------------------------

Система трёх уравнений на три неизвестных $(p_j, S_{g,j}, R_{s,j})$
(раздел~\ref{sssec:variable_switching}) решается по~абсолютно неявной
схеме методом Ньютона. Полный вектор неизвестных размера~$3N$:
\[
  \mathbf{U} = (p_1, S_{g,1}, R_{s,1},\;\;
                p_2, S_{g,2}, R_{s,2},\;\; \ldots,\;\;
                p_N, S_{g,N}, R_{s,N})^T.
\]
Матрица Якоби~--- блочно-трёхдиагональная с~блоками $3\times 3$.
Третья строка блока зависит от~режима ячейки (режим~1 или~2),
но~размерность системы всегда~$3N$.


\subsubsection{Потоки через грани}
 
Геометрическая трансмиссибильность грани~$j+\half$
в~логарифмической радиальной сетке:
\begin{equation}\label{eq:transmissibility}
  \mathcal{T}_{j+\half} = \frac{2\pi\,h\,k}{\ln(r_{j+1}/r_j)}.
\end{equation}
Делитель~--- $\ln(r_{j+1}/r_j)$, а~не~$\Delta r$ (следствие точного
решения для радиального потока в~кольцевом слое).
Абсолютная проницаемость~$k$ постоянна, поэтому интерполяция~$k$
на~грань не~требуется.

\paragraph{Потоки компонентов.}
Поток нефтяного компонента через грань:
\begin{equation}\label{eq:flux_oil}
  G_{j+\half}^{o} = \mathcal{T}_{j+\half}\,
    \left(\frac{k_{ro}}{\mu_o\,B_o}\right)_{\!j+\half}
    \!(p_{j+1} - p_j).
\end{equation}
Поток газового компонента:
\begin{equation}\label{eq:flux_gas}
  G_{j+\half}^{g}
  = \mathcal{T}_{j+\half}\!\left[
      \left(\frac{R_s\,k_{ro}}{\mu_o\,B_o}\right)_{\!j+\half}
      + \left(\frac{k_{rg}}{\mu_g\,B_g}\right)_{\!j+\half}
    \right](p_{j+1}-p_j).
\end{equation}

\paragraph{Upstream-взвешивание.}
Подвижности фаз $k_{r\alpha}/(\mu_\alpha\,B_\alpha)$ зависят от
насыщенности~$S_g$ (через ОФП) и~могут иметь разрыв на~фронте
дегазации. На~грани они берутся из~\textbf{upstream}-ячейки~---
той, откуда течёт фаза~$\alpha$:
\[
  \left(\frac{k_{r\alpha}}{\mu_\alpha\,B_\alpha}\right)_{\!j+\half}
  = \begin{cases}
      (k_{r\alpha}/\mu_\alpha B_\alpha)_j, & p_{j+1} > p_j
        \;\text{(поток влево)}, \\
      (k_{r\alpha}/\mu_\alpha B_\alpha)_{j+1}, & p_{j+1} < p_j
        \;\text{(поток вправо)}.
    \end{cases}
\]
Upstream обеспечивает монотонность схемы и~подавляет осцилляции.

\paragraph{Интерполяция скалярных свойств.}
Свойства, не~зависящие от~насыщенности (теплопроводность~$\leff$),
интерполируются на~грань \textbf{гармоническим средним}
с~логарифмическими весами:
\begin{equation}\label{eq:face_interp}
  A_{j+\half} = \frac{1}{\dfrac{\ln(r_{j+\half}/r_j)}{A_j}
    + \dfrac{\ln(r_{j+1}/r_{j+\half})}{A_{j+1}}}.
\end{equation}
Логарифмические веса корректны для радиальной геометрии
(аналогия с~последовательными сопротивлениями в~кольцевых слоях).


\subsubsection{Невязка}

Невязки формулируются через \textbf{исходные консервативные}
уравнения, что гарантирует строгое сохранение массы.

Для ячейки~$j$ интегрируем каждое уравнение по объёму~$V_j$
и~аппроксимируем временную производную разностью. Невязки
\textbf{одинаковы} для обоих режимов~--- различаются только
неизвестные и~способ вычисления $R_s$ и~$S_g$.

\paragraph{Невязка баланса нефтяного компонента:}
\begin{equation}\label{eq:residual_oil}
  \boxed{
    F_j^{\text{oil}} = V_j\,\phi\,
      \frac{(1 - S_{g,j}^{n+1})/B_{o,j}^{n+1}
            - (1 - S_{g,j}^n)/B_{o,j}^n}{\Delta t}
      + \left(G_{j+\half}^{o} - G_{j-\half}^{o}\right) = 0,
  }
\end{equation}
где $G_{j+\half}^{o} = \mathcal{T}_{j+\half}\,
(\Lambda_o/B_o)_{j+\half}\,(p_{j+1} - p_j)$~--- поток нефтяного
компонента через грань. Подвижность upstream по~направлению потока.

\paragraph{Невязка баланса газового компонента:}
\begin{equation}\label{eq:residual_gas}
  \boxed{
    F_j^{\text{gas}} = V_j\,\phi\,
      \frac{\left[\frac{R_{s,j}^{n+1}(1-S_{g,j}^{n+1})}{B_{o,j}^{n+1}}
            + \frac{S_{g,j}^{n+1}}{B_{g,j}^{n+1}}\right]
            - \left[\frac{R_{s,j}^n(1-S_{g,j}^n)}{B_{o,j}^n}
            + \frac{S_{g,j}^n}{B_{g,j}^n}\right]}{\Delta t}
      + \left(G_{j+\half}^{g} - G_{j-\half}^{g}\right) = 0,
  }
\end{equation}
где полный поток газового компонента через грань~$G_{j+\half}^{g}$
определён в~\eqref{eq:flux_gas}.

\paragraph{Невязка замыкающего уравнения:}

В~\textbf{режиме~1} (двухфазная зона, $S_g > 0$):
\begin{equation}\label{eq:residual_eq_r1}
  F_j^{\text{eq}} = R_{s,j} - R_s^{\text{eq}}(p_j) = 0.
\end{equation}

В~\textbf{режиме~2} (однофазная зона):
\begin{equation}\label{eq:residual_eq_r2}
  F_j^{\text{eq}} = S_{g,j} = 0.
\end{equation}

Итого для каждой ячейки~--- три невязки: $F_j^{\text{oil}}$,
$F_j^{\text{gas}}$, $F_j^{\text{eq}}$.

\paragraph{Режимы и невязки балансов масс.}
Невязки~$F^{\text{oil}}$~\eqref{eq:residual_oil}
и~$F^{\text{gas}}$~\eqref{eq:residual_gas} имеют одну и~ту~же форму
в~обоих режимах. Различия проявляются в~значениях~$S_g$ и~$R_s$:
\begin{itemize}[nosep]
  \item \textbf{Режим~1}: все три величины $(p, S_g, R_s)$ ненулевые;
        связь $R_s = R_s^{\text{eq}}(p)$ обеспечивается замыкающим
        уравнением~\eqref{eq:residual_eq_r1}.
  \item \textbf{Режим~2}: $S_g = 0$ (обеспечивается~\eqref{eq:residual_eq_r2}),
        $k_{rg} = 0$, $\Lambda_g = 0$. Величина~$R_s$ определяется
        из~$F^{\text{gas}} = 0$.
\end{itemize}


\subsubsection{Матрица Якоби}\label{sssec:jacobian}

Блочно-трёхдиагональная с~блоками $3\times 3$. Для ячейки~$j$
диагональный блок:
\begin{equation}\label{eq:jacobian_block}
  J_{j,j} = \begin{pmatrix}
    \dfrac{\partial F_j^{\text{oil}}}{\partial p_j} &
    \dfrac{\partial F_j^{\text{oil}}}{\partial S_{g,j}} &
    \dfrac{\partial F_j^{\text{oil}}}{\partial R_{s,j}} \\[10pt]
    \dfrac{\partial F_j^{\text{gas}}}{\partial p_j} &
    \dfrac{\partial F_j^{\text{gas}}}{\partial S_{g,j}} &
    \dfrac{\partial F_j^{\text{gas}}}{\partial R_{s,j}} \\[10pt]
    \dfrac{\partial F_j^{\text{eq}}}{\partial p_j} &
    \dfrac{\partial F_j^{\text{eq}}}{\partial S_{g,j}} &
    \dfrac{\partial F_j^{\text{eq}}}{\partial R_{s,j}}
  \end{pmatrix}.
\end{equation}

Внедиагональные блоки $J_{j,j\pm1}$~--- производные невязок ячейки~$j$
по~неизвестным соседних ячеек; содержат только производные потоков.
Третья строка внедиагональных блоков~--- нулевая (замыкающее уравнение
локально и~не~зависит от~соседей).

Первые две строки блока~--- производные балансов масс~--- \textbf{одинаковы}
для обоих режимов. Различается только \textbf{третья строка}.

\paragraph{Первые две строки (балансы масс).}

Производные $F^{\text{oil}}$ и~$F^{\text{gas}}$ по~всем трём
неизвестным $(p, S_g, R_s)$. Эти строки \textbf{одинаковы} для
обоих режимов.

Производные \textbf{члена накопления нефтяного компонента}
$(1-S_g)/B_o(p, R_s)$ (цепное правило, $B_o$ зависит от~$p$
и~$R_s$ независимо):
\begin{align}
  \frac{\partial}{\partial p}\left(\frac{1-S_g}{B_o}\right)
    &= (1-S_g)\,\frac{c_o}{B_o},
     \label{eq:daccum_oil_dp} \\[4pt]
  \frac{\partial}{\partial S_g}\left(\frac{1-S_g}{B_o}\right)
    &= -\frac{1}{B_o},
     \label{eq:daccum_oil_dSg} \\[4pt]
  \frac{\partial}{\partial R_s}\left(\frac{1-S_g}{B_o}\right)
    &= -(1-S_g)\,\frac{1}{B_o^2}\,\frac{\partial B_o}{\partial R_s},
     \label{eq:daccum_oil_dRs}
\end{align}
где $c_o = -(1/B_o)\,\partial B_o/\partial p|_{R_s}$~--- сжимаемость
при фиксированном составе.

Производные \textbf{члена накопления газового компонента}
$R_s(1-S_g)/B_o + S_g/B_g$:
\begin{align}
  \frac{\partial}{\partial p}\!\left(\frac{R_s(1\!-\!S_g)}{B_o}
    \!+\! \frac{S_g}{B_g}\right)
    &= \frac{R_s(1\!-\!S_g)\,c_o}{B_o}
       + \frac{S_g\,c_g}{B_g},
     \label{eq:daccum_gas_dp} \\[4pt]
  \frac{\partial}{\partial S_g}\!\left(\frac{R_s(1\!-\!S_g)}{B_o}
    \!+\! \frac{S_g}{B_g}\right)
    &= -\frac{R_s}{B_o} + \frac{1}{B_g},
     \label{eq:daccum_gas_dSg} \\[4pt]
  \frac{\partial}{\partial R_s}\!\left(\frac{R_s(1\!-\!S_g)}{B_o}
    \!+\! \frac{S_g}{B_g}\right)
    &= \frac{1\!-\!S_g}{B_o}
       - \frac{R_s(1\!-\!S_g)}{B_o^2}\,\frac{\partial B_o}{\partial R_s}.
     \label{eq:daccum_gas_dRs}
\end{align}

Производные \textbf{потоков} через грани по~$(p, S_g, R_s)$:
\begin{itemize}[nosep]
  \item по~$p$: перепад давления $+$ зависимость $\mu_o(p)$,
        $\mu_g(p)$, $B_o(p, R_s)$, $B_g(p)$;
  \item по~$S_g$: через ОФП ($dk_{ro}/dS_g$, $dk_{rg}/dS_g$);
  \item по~$R_s$: через $B_o(p, R_s)$, $\mu_o(p, R_s)$
        и~множитель~$R_s$ в~потоке газового компонента.
\end{itemize}


\paragraph{Третья строка (замыкающее уравнение).}

В~\textbf{режиме~1} (двухфазная зона), третье
уравнение~\eqref{eq:residual_eq_r1}:
$F^{\text{eq}} = R_s - R_s^{\text{eq}}(p) = 0$. Производные:
\begin{equation}\label{eq:deq_r1}
  \frac{\partial F^{\text{eq}}}{\partial p} =
    -\frac{\partial R_s^{\text{eq}}}{\partial p},
  \qquad
  \frac{\partial F^{\text{eq}}}{\partial S_g} = 0,
  \qquad
  \frac{\partial F^{\text{eq}}}{\partial R_s} = 1.
\end{equation}

В~\textbf{режиме~2} (однофазная зона), третье
уравнение~\eqref{eq:residual_eq_r2}:
$F^{\text{eq}} = S_g = 0$. Производные:
\begin{equation}\label{eq:deq_r2}
  \frac{\partial F^{\text{eq}}}{\partial p} = 0,
  \qquad
  \frac{\partial F^{\text{eq}}}{\partial S_g} = 1,
  \qquad
  \frac{\partial F^{\text{eq}}}{\partial R_s} = 0.
\end{equation}

Третья строка внедиагональных блоков~--- нулевая (замыкающее
уравнение не~зависит от~соседних ячеек).


\subsubsection{Переключение режимов}\label{sssec:switching_disc}

После каждой итерации Ньютона в~каждой ячейке проверяется:
\begin{enumerate}[nosep]
  \item Если ячейка в~\textbf{режиме~1} и~$S_{g,j}^{(k+1)} < 0$:
        переключить на~\textbf{режим~2}. Положить $S_{g,j} = 0$,
        заменить третью строку блока якобиана
        на~\eqref{eq:deq_r2}.
  \item Если ячейка в~\textbf{режиме~2}
        и~$R_{s,j}^{(k+1)} > R_s^{\text{eq}}(p_j^{(k+1)})$:
        переключить на~\textbf{режим~1}. Положить
        $R_{s,j} = R_s^{\text{eq}}(p_j)$, $S_{g,j}$~--- малое
        положительное значение. Заменить третью строку
        на~\eqref{eq:deq_r1}.
\end{enumerate}

Размерность системы не~меняется ($3N$). Структура блоков
$3\times 3$ сохраняется. Меняется только третья строка
диагонального блока.


\subsubsection{Итерационная процедура}

\begin{enumerate}[nosep]
  \item Начальное приближение:
        $\mathbf{U}^{(0)} = \mathbf{U}^n$ (включая режимы ячеек).
  \item Вычислить PVT и~ОФП при $\mathbf{U}^{(k)}$.
  \item Вычислить невязки $F^{\text{oil}}$, $F^{\text{gas}}$,
        $F^{\text{eq}}$ для каждой ячейки.
  \item Собрать блочно-трёхдиагональную матрицу Якоби~$J$
        (блоки $3\times 3$; третья строка зависит от~режима ячейки).
  \item Применить граничные условия (модификация коэффициентов).
  \item Решить СЛАУ: $J\,\delta\mathbf{U} = -\mathbf{F}$
        (блочная прогонка с~блоками $3\times 3$, сложность~$O(N)$).
  \item Обновить: $\mathbf{U}^{(k+1)} = \mathbf{U}^{(k)} + \delta\mathbf{U}$.
  \item Проверить переключение режимов
        (раздел~\ref{sssec:switching_disc}).
  \item Контроль физичности: $0 \leqslant S_g \leqslant 1 - S_{or}$,
        $0 \leqslant R_s \leqslant R_s^{\text{eq}}(p)$.
        Если нарушено~--- уменьшить~$\Delta t$, начать шаг заново.
  \item Проверить сходимость (раздел~\ref{ssec:convergence}).
\end{enumerate}

Блочная прогонка~--- обобщение скалярной прогонки Томаса на блочные
системы. Каждый шаг~--- обращение блока $3\times 3$. Сложность~$O(N)$.

Сходимость Ньютона квадратичная при хорошем начальном приближении
($p^n$); типично 2--4 итерации на~шаг.


%-----------------------------------------------------------------------------
\subsection{Уравнение энергии}\label{ssec:energy_disc}
%-----------------------------------------------------------------------------

Уравнение энергии~\eqref{eq:energy_linear} интегрируется
по~объёму ячейки~$V_j$. Все коэффициенты вычисляются
по~известным $p^{n+1}$, $u^{n+1}$, $S_g^{n+1}$ и~замораживаются
на~$T^n$. Неизвестная~--- $T_j^{n+1}$.

\paragraph{Upstream-оператор.}
Для произведения потока~$\mathcal{G}$ на~скалярную величину~$\psi$
на~грани~$j+\half$ используется аппроксимация против потока:
\begin{equation}\label{eq:upstream_operator}
  \langle\mathcal{G}\,\psi\rangle_{j+\half}
  = \mathcal{G}_{j+\half}^+\,\psi_j
  + \mathcal{G}_{j+\half}^-\,\psi_{j+1},
\end{equation}
где $\mathcal{G}^+ = \max(\mathcal{G},\,0)$,
$\mathcal{G}^- = \min(\mathcal{G},\,0)$.
Значение~$\psi$ берётся из~ячейки, откуда направлен
поток~$\mathcal{G}$.

\paragraph{Расход фазы через грань.}
Объёмный расход фазы~$\alpha$ (при пластовых условиях) через
грань~$j+\half$:
\begin{equation}\label{eq:phase_flux}
  \mathcal{Q}_{\alpha,j+\half}
  = \mathcal{T}_{j+\half}\,
    \left(\frac{k_{r\alpha}}{\mu_\alpha}\right)_{\!j+\half}
    \!(p_{j+1} - p_j),
\end{equation}
где подвижность~$k_{r\alpha}/\mu_\alpha$ берётся из~upstream-ячейки.

\paragraph{Конвективный поток энергии через грань.}
\begin{equation}\label{eq:energy_conv_face}
  \mathcal{F}_{j+\half}
  = \sum_\alpha
    \langle\Phi_\alpha\,T\rangle_{j+\half},
\end{equation}
где $\Phi_\alpha = \rho_\alpha\,c_{p,\alpha}\,
\mathcal{Q}_\alpha$~--- конвективный коэффициент переноса
фазы~$\alpha$ через грань [Вт/К].
Upstream определяется \textbf{раздельно}
по~каждой фазе (фазы могут течь в~разные стороны).

\paragraph{Кондуктивный поток через грань.}
\begin{equation}\label{eq:heat_cond_disc}
  Q_{\lambda,j+\half} = \frac{2\pi\,h}{\ln(r_{j+1}/r_j)}\,
    {\leff}^{j+\half}\,(T_{j+1} - T_j),
\end{equation}
где~${\leff}^{j+\half}$~--- гармоническое среднее
с~логарифмическими весами (ср.~\eqref{eq:face_interp}):
\[
  {\leff}^{j+\half} =\left( \dfrac{\ln(r_{j+\half}/r_j)}{{\leff}_j}
    + \dfrac{\ln(r_{j+1}/r_{j+\half})}{{\leff}_{j+1}}\right)^{-1}.
\]

\paragraph{Эффект Джоуля--Томсона (пространственный).}
JT-поток через грань~$j+\half$:
\begin{equation}\label{eq:QJT_face}
  \mathcal{J}_{j+\half}
  = \sum_\alpha \mu_{\mathrm{JT},\alpha}^{j+\half}\,
    \langle\Phi_\alpha\,p\rangle_{j+\half},
\end{equation}
где~$\mu_{\mathrm{JT},\alpha}^{j+\half}$~--- коэффициент
Джоуля--Томсона на~грани (upstream по~фазе~$\alpha$).
Вклад JT в~уравнение ячейки~$j$~--- разность потоков:
\begin{equation}\label{eq:QJT_disc}
  \mathcal{J}_j = \mathcal{J}_{j+\half} - \mathcal{J}_{j-\half}.
\end{equation}

\paragraph{Адиабатический эффект (временной).}
Вычисляется поточечно в~центре ячейки (при~$c_r = 0$,
через~$\mu_{\mathrm{JT}}$~\eqref{eq:main_Qad_muJT}):
\begin{equation}\label{eq:Qad_disc}
  V_j\,{\Qad}_j = V_j\,\phi\,
    \sum_\alpha S_{\alpha,j}\,
    \bigl(1 + \rho_{\alpha,j}\,c_{p,\alpha,j}\,
    \mu_{\mathrm{JT},\alpha,j}\bigr)\,
    \frac{p_j^{n+1}-p_j^n}{\Delta t}.
\end{equation}

\paragraph{Теплота дегазации.}
Вычисляется поточечно:
\begin{equation}\label{eq:Qdeg_disc}
  V_j\,{\Qdeg}_j = -V_j\,L_j\,\mdot_j.
\end{equation}

Массовый темп дегазации~\eqref{eq:main_mdot} в~дискретной форме.
Временная часть~--- разность по~времени в~центре ячейки;
пространственная часть $(u_o/B_o)\,\partial R_s/\partial r$~---
upstream через потоки на~гранях.
Конвективный поток газосодержания через грань~$j+\half$:
\begin{equation}\label{eq:mdot_conv_face}
  \mathcal{R}_{j+\half}
  = \langle G^o\,R_s\rangle_{j+\half},
\end{equation}
где $G^o = G^o_{j+\half}$~--- поток нефтяного
\textbf{компонента} через грань~\eqref{eq:flux_oil}
[м$^3$ст/с] (не путать с~расходом фазы~$\mathcal{Q}_o$
[м$^3$пл/с]).
Массовый темп дегазации для ячейки~$j$:
\begin{equation}\label{eq:mdot_disc}
  V_j\,\mdot_j = -\rho_g^{\text{ст}}\!\left[
    V_j\,\phi\,\frac{1-S_{g,j}^{n+1}}{B_{o,j}^{n+1}}\,
      \frac{R_{s,j}^{n+1} - R_{s,j}^n}{\Delta t}
    + \bigl(\mathcal{R}_{j+\half} - \mathcal{R}_{j-\half}\bigr)
  \right].
\end{equation}
Формула единая для обоих режимов. В~режиме~2
($R_s = R_s^0 = \mathrm{const}$) обе части обращаются
в~нуль и~$\mdot_j = 0$ автоматически.

Все величины вычислены при~$p^{n+1}$, $S_g^{n+1}$ и~$T^n$.

\paragraph{Дискретное уравнение для ячейки~$j$.}
Интеграл уравнения~\eqref{eq:energy_linear} по~$V_j$:
\begin{equation}\label{eq:energy_disc_full}
  \boxed{
    V_j\,{\rhocpeff}_j\,\frac{T_j^{n+1} - T_j^n}{\Delta t}
    + \bigl(\mathcal{F}_{j+\half} - \mathcal{F}_{j-\half}\bigr)
    = \bigl(Q_{\lambda,j+\half} - Q_{\lambda,j-\half}\bigr)
    + \mathcal{J}_j
    + V_j\,({\Qad}_j + {\Qdeg}_j).
  }
\end{equation}
Конвективные потоки~$\mathcal{F}$ и~кондуктивные
потоки~$Q_\lambda$ вычисляются через грани ячейки (левую
$j{-}\half$ и~правую $j{+}\half$); JT-вклад $\mathcal{J}_j$
использует потоки на~обеих гранях~\eqref{eq:QJT_disc};
$\Qad$ и~$\Qdeg$~--- поточечно в~центре.

\paragraph{Трёхдиагональная СЛАУ.}
Подставляя~\eqref{eq:energy_conv_face} и~\eqref{eq:heat_cond_disc}
в~\eqref{eq:energy_disc_full}, собираем слагаемые
при~$T_{j-1}^{n+1}$, $T_j^{n+1}$, $T_{j+1}^{n+1}$:
\begin{equation}\label{eq:energy_tridiag}
  a_j\,T_{j-1}^{n+1} + b_j\,T_j^{n+1} + c_j\,T_{j+1}^{n+1} = d_j,
\end{equation}
где
\begin{align}
  a_j &= -Q_{\lambda,j-\half}^T - \Phi_{j-\half}^{+},
    \label{eq:tridiag_a} \\
  c_j &= -Q_{\lambda,j+\half}^T + \Phi_{j+\half}^{-},
    \label{eq:tridiag_c} \\
  b_j &= \frac{V_j\,{\rhocpeff}_j}{\Delta t}
    + Q_{\lambda,j-\half}^T + Q_{\lambda,j+\half}^T
    + \Phi_{j+\half}^{+} - \Phi_{j-\half}^{-},
    \label{eq:tridiag_b} \\
  d_j &= \frac{V_j\,{\rhocpeff}_j}{\Delta t}\,T_j^n
    + \mathcal{J}_j
    + V_j\,({\Qad}_j + {\Qdeg}_j).
    \label{eq:tridiag_d}
\end{align}
Здесь $Q_{\lambda,j\pm\half}^T = 2\pi h\,{\leff}^{j\pm\half}/
\ln(r_{j\pm 1}/r_j)$~--- кондуктивная трансмиссибильность ($>0$);
$\Phi_{j+\half}^{+} = \sum_\alpha \Phi_{\alpha,j+\half}^+ \geqslant 0$,
$\Phi_{j+\half}^{-} = \sum_\alpha \Phi_{\alpha,j+\half}^- \leqslant 0$~---
суммарные upstream-разделённые потоки~\eqref{eq:upstream_operator};
$\mathcal{J}_j$~--- JT-вклад~\eqref{eq:QJT_disc} (известное число).
Коэффициенты $a_j \leqslant 0$, $c_j \leqslant 0$, $b_j > 0$~---
диагональное преобладание обеспечено.

СЛАУ~\eqref{eq:energy_tridiag} решается скалярной прогонкой
(алгоритм Томаса) без итераций.


%-----------------------------------------------------------------------------
\subsection{Граничные условия (дискретная форма)}\label{ssec:bc_disc}
%-----------------------------------------------------------------------------

Граничные условия реализуются через \textbf{потоки
на~граничных гранях}. Давление на~границе задано, поток
через граничную грань вычисляется как разность давлений
между центром граничной ячейки и~границей. Все три
неизвестных $(p_j, S_{g,j}, R_{s,j})$ в~граничных ячейках~---
полноценные переменные, определяемые из~балансов масс.

Трансмиссибильность граничной грани определяется по~общей
формуле~\eqref{eq:transmissibility}, где вместо расстояния
между центрами соседних ячеек используется расстояние от~центра
граничной ячейки до границы области:
\[
  \mathcal{T}_{\half} = \frac{2\pi\,h\,k}{\ln(r_0/\rw)},
  \qquad
  \mathcal{T}_{N+\half} = \frac{2\pi\,h\,k}{\ln(\Rext/r_{N-1})}.
\]

\paragraph{Забой ($r = \rw$, левая грань ячейки $j = 0$).}
Давление на~грани $r_{\half} = \rw$ задано: $p\big|_{r_w} = \pwf$.
Потоки компонентов через левую грань:
\begin{align}
  G_{\half}^{o} &= \mathcal{T}_{\half}\,
    \left(\frac{k_{ro}}{\mu_o\,B_o}\right)_{\!0}
    \,(p_0 - \pwf),
    \label{eq:bc_flux_well} \\
  G_{\half}^{g} &= \mathcal{T}_{\half}\!\left[
    \left(\frac{R_s\,k_{ro}}{\mu_o\,B_o}\right)_{\!0}
    + \left(\frac{k_{rg}}{\mu_g\,B_g}\right)_{\!0}
    \right](p_0 - \pwf).
    \label{eq:bc_flux_well_gas}
\end{align}
Подвижности берутся из~ячейки~$j = 0$
(upstream: поток направлен к~скважине).

Потоки $G_{\half}^{o}$, $G_{\half}^{g}$ входят в~невязки
ячейки~$j = 0$ наравне с~потоками через правую грань~$G_{3/2}$.
Давление~$p_0$~--- \textbf{неизвестная}, определяемая
из~баланса масс; $\pwf$~--- заданный параметр.
Дебит скважины (при ст.~условиях) определяется
потоком~$G_{\half}^o$.

\paragraph{Контур питания ($r = \Rext$, правая грань ячейки
$j = N{-}1$).}
Давление на~грани $r_{N+\half} = \Rext$ задано:
$p\big|_{\Rext} = \ppl$. На~контуре нефть однофазная
($p = \ppl > \pb$): $S_g\big|_{\Rext} = 0$,
$k_{rg} = 0$, $R_s\big|_{\Rext} = R_s^0$. Потоки:
\begin{align}
  G_{N+\half}^{o} &= \mathcal{T}_{N+\half}\,
    \left(\frac{k_{ro}}{\mu_o\,B_o}\right)_{\!N-1}
    \,(\ppl - p_{N-1}),
    \label{eq:bc_flux_ext} \\
  G_{N+\half}^{g} &= \mathcal{T}_{N+\half}\,
    \frac{R_s^0\,(k_{ro})_{N-1}}{(\mu_o\,B_o)_{N-1}}
    \,(\ppl - p_{N-1}).
    \label{eq:bc_flux_ext_gas}
\end{align}
Свободного газа нет ($k_{rg}=0$); поток газового компонента
обусловлен только растворённым газом с~составом~$R_s^0$.

\paragraph{Источниковые члены на~граничных гранях.}
JT-поток~\eqref{eq:QJT_face} и~конвективный поток
газосодержания~\eqref{eq:mdot_conv_face} на~граничных гранях
вычисляются по~тем же формулам, что и~на~внутренних, с~заменой:
\begin{itemize}[nosep]
  \item на~забое: $p_{j+1} \to \pwf$, $R_{s,j+1} \to R_{s,0}$
    в~upstream-операторе~\eqref{eq:upstream_operator}
    (давление на~грани задано, $R_s$~--- из~ячейки~$j = 0$);
  \item на~контуре: $p_{j+1} \to \ppl$,
    $R_{s,j+1} \to R_s^0$
    (приток с~контура имеет известный состав).
\end{itemize}
Расход фазы через граничную грань~\eqref{eq:phase_flux}
вычисляется с~граничной трансмиссибильностью $\mathcal{T}_{\half}$
или~$\mathcal{T}_{N+\half}$. Далее
$\Phi_{\alpha} = \rho_\alpha\,c_{p,\alpha}\,\mathcal{Q}_\alpha$
определяется по~общей формуле.

\paragraph{Температура.}
\begin{itemize}[nosep]
  \item $r = \Rext$: условие I~рода, $T\big|_{\Rext} = \Tpl$.
    Конвективный поток энергии через правую грань
    ячейки~$j = N{-}1$: upstream берёт~$T = \Tpl$
    (приток с~контура).
    Кондуктивный поток:
    $Q_{\lambda,N+\half}
    = \dfrac{2\pi\,h}{\ln(\Rext/r_{N-1})}\,
    {\leff}_{N-1}\,(\Tpl - T_{N-1})$
    (${\leff}$ берётся из~ячейки~$j = N{-}1$, так как
    на~границе нет соседней ячейки для интерполяции).
  \item $r = \rw$: $\partial T/\partial r = 0$.
    Кондуктивный поток через левую грань обнуляется:
    $Q_{\lambda,\half} = 0$.
    Конвективный поток: флюид уходит в~скважину, upstream
    берёт~$T = T_0^{n+1}$ из~ячейки~$j = 0$~---
    дополнительного граничного условия не~требуется.
\end{itemize}


%=============================================================================
\section{Интегрирование по времени}\label{sec:time}
%=============================================================================

%-----------------------------------------------------------------------------
\subsection{Схема по времени}\label{ssec:time_scheme}
%-----------------------------------------------------------------------------

Все уравнения решаются по~\textbf{абсолютно неявной} схеме.
Конвективные члены аппроксимируются \textbf{схемой против потока}
(upwind). Схема имеет первый порядок точности по~времени
и~пространству, но~безусловно устойчива и~консервативна~--- это
ключевое требование.

\begin{itemize}[nosep]
  \item \textbf{Уравнения давления и~насыщенности}~--- совместная
    полностью неявная система, метод Ньютона. Неизвестные
    $(p, S_g, R_s)$ определяются одновременно на~уровне~$n{+}1$.
    PVT и~ОФП обновляются на~каждой итерации.
  \item \textbf{Уравнение энергии}~--- неявная по~$T$. Коэффициенты
    вычислены по~$p^{n+1}$, $u^{n+1}$, $S_g^{n+1}$ и~заморожены.
    СЛАУ линейна на~каждом шаге.
\end{itemize}

Устойчивость и~консервативность важнее порядка для задач
с~разгазированием (резкое изменение свойств на фронте). Повышение
порядка (Crank--Nicolson, TVD) может быть добавлено позже.


%-----------------------------------------------------------------------------
\subsection{Алгоритм одного шага}\label{ssec:algorithm}
%-----------------------------------------------------------------------------

\begin{enumerate}[leftmargin=*, label=\textbf{Шаг \arabic*.}]

\item \textbf{Определить $\pwf$.}
  Найти интервал истории, содержащий~$t^n$: $\pwf = \pwf^{(k)}$.

\item \textbf{Итерации Ньютона для $(p, X)$, $X = S_g$ или~$R_s$.}
  Начальное приближение: $\mathbf{U}^{(0)} = \mathbf{U}^n$
  (включая режимы ячеек с~предыдущего шага).

  Цикл ($k = 0, 1, \ldots$):
  \begin{enumerate}[nosep, label=(\alph*)]
    \item Вычислить PVT$(p^{(k)}, \Tpl)$ и~их производные.
    \item Вычислить ОФП и~их производные (в~режиме~2: $k_{rg}=0$).
    \item Вычислить невязки $F^{\text{oil}}$, $F^{\text{gas}}$
          для каждой ячейки.
    \item Собрать блочно-трёхдиагональную матрицу Якоби~$J$
          (блоки $J^{[1]}$ или~$J^{[2]}$ по~режиму ячейки,
          раздел~\ref{sssec:jacobian}).
    \item Применить граничные условия (модификация коэффициентов).
    \item Решить СЛАУ: $J\,\delta\mathbf{U} = -\mathbf{F}$
          (блочная прогонка).
    \item Обновить: $\mathbf{U}^{(k+1)} = \mathbf{U}^{(k)}
          + \delta\mathbf{U}$.
    \item Проверить переключение режимов
          (раздел~\ref{sssec:switching_disc}).
    \item Контроль физичности. Если нарушено~--- уменьшить~$\Delta t$,
          начать заново.
    \item Проверить сходимость (раздел~\ref{ssec:convergence}).
          Если выполнено~--- выход.
  \end{enumerate}
  Результат: $p^{n+1}$, $S_g^{n+1}$ (или~$R_s^{n+1}$ в~однофазных
  ячейках).

\item \textbf{Вычислить скорости Дарси.}
  $u_o^{n+1}$, $u_g^{n+1}$ из~$p^{n+1}$.

\item \textbf{Решение тепловой задачи.}
  \begin{enumerate}[nosep, label=(\alph*)]
    \item Вычислить PVT$(p^{n+1}, T^n)$: $\rho_o$, $\rho_g$, $c_{p,o}$,
          $c_{p,g}$, $\alpha_o$, $\alpha_g$, $L$, $\leff$.
    \item Вычислить коэффициенты: $\rhocpeff$, $\rhocpconv$, $\QJT$,
          $\Qad$.
    \item Вычислить $\Qdeg$: $\Qdeg < 0$~--- неявно (в~матрицу),
          $\Qdeg > 0$~--- явно (в~правую часть).
    \item Собрать трёхдиагональную матрицу уравнения энергии.
    \item Применить граничные условия.
    \item Решить СЛАУ (прогонка) $\to$ $T^{n+1}$.
  \end{enumerate}

\item \textbf{Контроль шага по~времени.}
  Адаптация $\Delta t$ в~документе не~описывается~--- стратегия
  определяется при реализации в~коде.

\end{enumerate}


%-----------------------------------------------------------------------------
\subsection{Критерии сходимости}\label{ssec:convergence}
%-----------------------------------------------------------------------------

\paragraph{CFL-ограничение отсутствует:}
схема полностью неявная, безусловно устойчива. Шаг по~времени
ограничивается только точностью, не~устойчивостью.

\paragraph{Сходимость Ньютона.}
Критерий~--- \textbf{покомпонентный}, для каждой ячейки~$j$
должно выполняться хотя бы одно из двух условий:

\emph{Абсолютное:}
\[
  |F_j^{\text{oil}}| < \varepsilon_{\text{oil}}^{\text{abs}}
  \quad\text{и}\quad
  |F_j^{\text{gas}}| < \varepsilon_{\text{gas}}^{\text{abs}}.
\]

\emph{Относительное:}
\[
  \frac{|F_j^{\text{oil}}|}{|F_{j,0}^{\text{oil}}|}
    < \varepsilon_{\text{oil}}^{\text{rel}}
  \quad\text{и}\quad
  \frac{|F_j^{\text{gas}}|}{|F_{j,0}^{\text{gas}}|}
    < \varepsilon_{\text{gas}}^{\text{rel}},
\]
где $F_{j,0}$~--- невязка на~первой итерации.

Пределы погрешности задаются пользователем
в~конфигурационном файле.

\paragraph{Контроль физичности:}
\begin{itemize}[nosep]
  \item $0 \leqslant S_g \leqslant 1 - S_{or}$ (в~режиме~1);
  \item $0 \leqslant R_s \leqslant R_s^{\text{eq}}(p)$ (в~режиме~2);
  \item $p > 0$;
  \item $T > 0$.
\end{itemize}
Нарушение любого~--- $\Delta t$ уменьшается, шаг повторяется.


%=============================================================================
\section*{Переходные режимы}
\addcontentsline{toc}{section}{Переходные режимы}
%=============================================================================

Переключение неизвестных между режимами~1 и~2 описано
в~разделах~\ref{sssec:variable_switching} (физика)
и~\ref{sssec:switching_disc} (дискретная реализация).

При $S_g < S_{gc}$ (критическая газонасыщенность): $k_{rg} = 0$
тождественно, газ неподвижен~--- газовая фаза присутствует, но
не~фильтруется.

\paragraph{Ограничение: исчезновение нефтяной фазы.}
Текущая постановка не~предусматривает ситуацию $S_g = 1$
(полное отсутствие нефтяной фазы). При $S_o = 0$ баланс нефтяного
компонента~\eqref{eq:mass_oil} вырождается ($0 = 0$), якобиан
теряет ранг, и~метод Ньютона не~сходится.

Это ограничение не~является существенным для рассматриваемой задачи
по~двум причинам:
\begin{enumerate}[nosep]
  \item Нефть непрерывно поступает через контур питания ($S_g = 0$
        при $r = \Rext$), поэтому нефтяная фаза присутствует
        во~всей расчётной области.
  \item При $S_{or} > 0$ (остаточная нефтенасыщенность) модель ОФП
        ограничивает газонасыщенность: $S_g \leqslant 1 - S_{or}$.
        Нефтяная фаза не~может исчезнуть полностью~--- остаточная
        нефть неподвижна, но~присутствует.
\end{enumerate}
При необходимости моделирования полного вытеснения нефти газом
потребуется введение третьего режима (аналогичного режиму~2, но
для~газовой однофазной зоны).


%=============================================================================
\appendix

%=============================================================================
\section{Вывод уравнения для темпа выделения свободного газа}\label{app:degas_rate}
%=============================================================================

Ниже приведён вывод уравнения баланса свободного газа
и~определения массового темпа дегазации~$\mdot$.

Уравнение баланса нефтяного компонента~\eqref{eq:mass_oil} записано
в~объёмных единицах нефтяного компонента при стандартных условиях
[м$^3$ст\,нефти/(м$^3\cdot$с)], плотность~$\rho_o^{\text{ст}}$
была сокращена из~обеих частей. Аналогично,
уравнение~\eqref{eq:mass_gas} записано в~объёмных единицах газового
компонента при стандартных условиях [м$^3$ст\,газа/(м$^3\cdot$с)],
плотность~$\rho_g^{\text{ст}}$ сокращена.

Газосодержание $R_s$ имеет размерность
[м$^3$ст\,газа\,/\,м$^3$ст\,нефти]. Поэтому результат умножения
$R_s$ на уравнение~\eqref{eq:mass_oil} имеет ту же размерность, что
и~\eqref{eq:mass_gas}~--- м$^3$ст\,газа/(м$^3\cdot$с). Это
допускает вычитание одного уравнения из~другого:
\begin{multline}\label{eq:gas_minus_Rs_oil}
  \pdt{}\!\left(\phi\!\left[
    \frac{R_s\,S_o}{B_o} + \frac{S_g}{B_g}
  \right]\right)
  - R_s\,\pdt{}\!\left(\phi\,\frac{S_o}{B_o}\right)
  \\
  + \divr{\left[
    \frac{R_s\,u_o}{B_o} + \frac{u_g}{B_g}
  \right]}
  - R_s\,\divr{\frac{u_o}{B_o}} = 0.
\end{multline}

Перегруппируем слагаемые, чтобы привести уравнение к~виду
баланса массы газа в~газовой фазе и~явно выделить вклад,
определяющий производство свободного газа:
\begin{equation}\label{eq:free_gas_rate}
  \boxed{
    \pdt{}\!\left(\phi\,\frac{S_g}{B_g}\right)
    + \divr{\frac{u_g}{B_g}}
    = -\phi\,\xio\,\pdt{R_s}
      - \frac{u_o}{B_o}\,\pdr{R_s}.
  }
\end{equation}
Здесь слева~--- локальное накопление и~конвективный
перенос свободного газа, справа~--- источник газа за~счёт
разгазирования нефти.
Временная и~конвективная части разности соответственно:
\[
  \pdt{}\!\left(\phi\,R_s\,\xio\right)
  - R_s\,\pdt{}\!\left(\phi\,\xio\right)
  = \phi\,\xio\,\pdt{R_s}, \quad \xio = S_o/B_o.
\]
\[
  \divr{\frac{R_s\,u_o}{B_o}}
  - R_s\,\divr{\frac{u_o}{B_o}}
  = \frac{u_o}{B_o}\,\pdr{R_s}.
\]
Оба слагаемых описывают переход газового компонента из~растворённого
состояния в~свободное: первое~--- за~счёт локального изменения~$R_s$
во~времени, второе~--- за~счёт притока нефти с~иным газосодержанием.
Перенося их в~правую часть~\eqref{eq:free_gas_rate}, получаем
источник газовой фазы.

Правая часть~\eqref{eq:free_gas_rate} имеет размерность
м$^3$ст\,газа/(м$^3\cdot$с)~--- объёмный (при стандартных
условиях) темп появления свободного газа. Для перехода к~массовому
темпу~$\mdot$ [кг/(м$^3\cdot$с)] умножаем на~плотность газового
компонента при стандартных условиях~$\rho_g^{\text{ст}}$:
\begin{equation}\label{eq:mdot}
  \boxed{
    \mdot = -\rho_g^{\text{ст}}\!\left[
      \phi\,\xio\,\pdt{R_s}
      + \frac{u_o}{B_o}\,\pdr{R_s}
    \right].
  }
\end{equation}

При~$\mdot > 0$ газ выделяется из нефти; при~$\mdot < 0$~---
растворяется обратно.


%=============================================================================
\section{Вывод уравнения энергии}\label{app:energy_derivation}
%=============================================================================

Ниже приведён полный вывод уравнения
энергии~\eqref{eq:energy_final} из~первого начала термодинамики.

Рассмотрим неподвижный элемент пористой среды. Первое начало
термодинамики для открытой системы:
\begin{equation}\label{eq:first_law_open}
  \frac{dE_{\mathrm{total}}}{dt}
  = \dot{Q} - \dot{W}_{\text{внешн}}
    + \sum_{\text{входы}} \dot{m}_{\mathrm{in}}\,h_{\mathrm{in}}
    - \sum_{\text{выходы}} \dot{m}_{\mathrm{out}}\,h_{\mathrm{out}}.
\end{equation}

Определим входящие в~\eqref{eq:first_law_open} величины:
\begin{itemize}[nosep]
  \item $E_{\mathrm{total}}$~--- полная внутренняя энергия вещества
        внутри элемента (порода + флюиды), Дж/м$^3$;
  \item $\dot{Q}$~--- подвод тепла к~элементу за~счёт теплопроводности
        (кондукция через границы элемента), Вт/м$^3$;
  \item $\dot{W}_{\text{внешн}}$~--- мощность работы, совершаемой
        системой над~окружением, помимо работы проталкивания
        (в~рассматриваемой задаче $\dot{W}_{\text{внешн}} = 0$:
        нет подвижных стенок и~внешних механических воздействий),
        Вт/м$^3$;
  \item $\dot{m}_{\mathrm{in}}$, $\dot{m}_{\mathrm{out}}$~--- массовые
        расходы флюида, втекающего в~элемент и~вытекающего из~него
        через его границы, кг/(м$^3\cdot$с);
  \item $h_{\mathrm{in}}$, $h_{\mathrm{out}}$~--- удельные энтальпии
        втекающего и~вытекающего флюида, Дж/кг.
\end{itemize}

Для рассматриваемой системы
уравнение~\eqref{eq:first_law_open} в \textbf{гомогенном приближении} запишется следующим образом
\begin{equation}\label{eq:energy_cons}
  \pdt{E}
  + \divr{F}
  = \divr{\leff\,\pdr{T}},
\end{equation}
где
\begin{equation}\label{eq:E_expanded}
  E = (1 - \phi)\,(\rho c_p)_s\,T + \phi\sum_\alpha S_\alpha\,\rho_\alpha\,e_\alpha,
\end{equation}
\begin{equation}\label{eq:E_enthalpy_relation}
\phi\sum_\alpha S_\alpha\,\rho_\alpha\,e_\alpha = \phi\sum_\alpha S_\alpha\,\rho_\alpha\,h_\alpha
      - \phi\,p,
\end{equation}
\begin{equation}\label{eq:F_total}
  F = \sum_\alpha \rho_\alpha\,h_\alpha\,u_\alpha, \quad  h_\alpha = e_\alpha + p/\rho_\alpha.
\end{equation}

Определим энтальпию многокомпонентной нефтяной фазы, которая является смесью двух компонентов.
Её энтальпия разлагается через парциальные энтальпии:
\begin{equation}\label{eq:oil_enthalpy_components}
  \phi\,S_o\,\rho_o\,h_o
  = \phi\,\frac{S_o}{B_o}\!\left(
      \rho_o^{\text{ст}}\,\bar{h}_{oo}
      + R_s\,\rho_g^{\text{ст}}\,\bar{h}_{go}
    \right),
\end{equation}
где~$\bar{h}_{oo}(p, T, R_s)$~--- парциальная удельная
энтальпия\footnote{Парциальная удельная энтальпия компонента~$k$
в~фазе~--- это производная энтальпии фазы~$H$ по~массе~$m_k$ этого
компонента при постоянных давлении~$p$, температуре~$T$ и~массах
остальных компонентов~$m_{j\neq k}$:
$\bar{h}_k = (\partial H/\partial m_k)_{p,\,T,\,m_{j\neq k}}$.
Она характеризует вклад единицы массы данного компонента в~энтальпию
раствора и~отличается как от~энтальпии чистого вещества, так и~от~средней
энтальпии смеси.} нефтяного компонента в~нефтяной фазе,
$\bar{h}_{go}(p, T, R_s)$~--- парциальная удельная энтальпия газового
компонента в~нефтяной фазе. Обе величины зависят от~давления,
температуры и~фактического состава нефтяной фазы (газосодержания~$R_s$).
Плотность нефтяной фазы выразится следующим образом через плотности компонентов в~стандартных условиях
\begin{equation}\label{eq:rho_oil}
  \rho_o = \frac{\rho_o^{\text{ст}} + R_s\,\rho_g^{\text{ст}}}{B_o}.
\end{equation}
Газовая фаза содержит только один компонент (газ), поэтому её энтальпия
просто~$h_g(p,T)$.

В уравнении~\eqref{eq:energy_cons} выделим явно слагаемые, которые определяют
вклад в~общий баланс энергии и~отвечают (i)~фазовому переходу (дегазации),
(ii)~эффекту Джоуля--Томсона и (iii)~дроссельному эффекту.

В выражении полной энтальпии системы выделим явно части,
которые отвечают отдельным компонентам, а~не фазам.
Для внутренней энергии и~конвективного потока получим соответственно
\begin{equation}\label{eq:H_fluids}
  \phi\sum_\alpha S_\alpha\rho_\alpha h_\alpha
  = \underbrace{\phi\,\frac{S_o}{B_o}\,\rho_o^{\text{ст}}\,
      \bar{h}_{oo}}_{\text{нефтяной комп.}}
    + \underbrace{\phi\,\rho_g^{\text{ст}}\!\left[
        \frac{R_s\,S_o}{B_o}\,\bar{h}_{go}
        + \frac{S_g}{B_g}\,h_g
      \right]}_{\text{газовый комп.}}.
\end{equation}
\begin{equation}\label{eq:F_components}
  F = \underbrace{\frac{u_o}{B_o}\,\rho_o^{\text{ст}}\,
        \bar{h}_{oo}}_{\text{нефт.\ комп.}}
      + \underbrace{\rho_g^{\text{ст}}\!\left[
          \frac{R_s\,u_o}{B_o}\,\bar{h}_{go}
          + \frac{u_g}{B_g}\,h_g
        \right]}_{\text{газ.\ комп.}}.
\end{equation}

Вклад газового компонента в уравнении~\eqref{eq:energy_cons} содержит
два слагаемых~--- растворённый газ (с~$\bar{h}_{go}$) и~свободный газ
(с~$h_g$). Рассмотрим их по~отдельности. Из~\eqref{eq:H_fluids} и~\eqref{eq:F_components}
соберём все члены, содержащие~$\bar{h}_{go}$, и~раскроем производные

\begin{multline}\label{eq:F_components_der}
  \pdt{}\!\left(\phi\,\rho_g^{\text{ст}}\,
    \frac{R_s\,S_o}{B_o}\,\bar{h}_{go}\right)
  + \divr{\rho_g^{\text{ст}}\,
    \frac{R_s\,u_o}{B_o}\,\bar{h}_{go}} =
    \\
    \rho_g^{\text{ст}}\,\bar{h}_{go}\!\left[
    \pdt{}\!\left(\phi\,
      \frac{R_s\,S_o}{B_o}\right)
    + \divr{
      \frac{R_s\,u_o}{B_o}}
  \right]
  +
  \\
    \left(\phi\,\rho_g^{\text{ст}}\,
      \frac{R_s\,S_o}{B_o}\right)
    \pdt{\bar{h}_{go}}\!
    + \left(\rho_g^{\text{ст}}\,
      \frac{R_s\,u_o}{B_o}\right)\pdr{\bar{h}_{go} }.
\end{multline}
Вычитая уравнение баланса свободного
газа~\eqref{eq:free_gas_rate} из~уравнения баланса газового
компонента~\eqref{eq:mass_gas}, выразим сумму в квадратных скобках
через массовый темп производства свободного газа
\[
  \pdt{}\!\left(\phi\,\frac{R_s\,S_o}{B_o}\right)
  + \divr{\frac{R_s\,u_o}{B_o}}
  = \phi\,\xio\,\pdt{R_s}
    + \frac{u_o}{B_o}\,\pdr{R_s}
  = -\frac{\mdot}{\rho_g^{\text{ст}}},
\]
где последнее равенство следует из~определения~\eqref{eq:mdot}.
Следовательно, уравнение~\eqref{eq:F_components_der} примет следующий вид
\begin{multline}\label{eq:F_components_derived}
  \pdt{}\!\left(\phi\,\rho_g^{\text{ст}}\,
    \frac{R_s\,S_o}{B_o}\,\bar{h}_{go}\right)
  + \divr{\rho_g^{\text{ст}}\,
    \frac{R_s\,u_o}{B_o}\,\bar{h}_{go}} =
    \\
    -\mdot\bar{h}_{go}
  +
    \left(\phi\,\rho_g^{\text{ст}}\,
      \frac{R_s\,S_o}{B_o}\right)
    \pdt{\bar{h}_{go}}\!
    + \left(\rho_g^{\text{ст}}\,
      \frac{R_s\,u_o}{B_o}\right)\pdr{\bar{h}_{go} }
\end{multline}

Аналогично, соберём члены с~$h_g$ из~\eqref{eq:H_fluids},
\eqref{eq:F_components} и~раскроем производные:
\begin{multline}\label{eq:hg_components_der}
  \pdt{}\!\left(\phi\,\rho_g^{\text{ст}}\,
    \frac{S_g}{B_g}\,h_g\right)
  + \divr{\rho_g^{\text{ст}}\,
    \frac{u_g}{B_g}\,h_g} =
    \\
    \rho_g^{\text{ст}}\,h_g\!\left[
    \pdt{}\!\left(\phi\,
      \frac{S_g}{B_g}\right)
    + \divr{\frac{u_g}{B_g}}
  \right]
  +
  \\
    \left(\phi\,\rho_g^{\text{ст}}\,
      \frac{S_g}{B_g}\right)
    \pdt{h_g}\!
    + \left(\rho_g^{\text{ст}}\,
      \frac{u_g}{B_g}\right)\pdr{h_g}.
\end{multline}
Выражение в~квадратных скобках~--- левая часть баланса свободного
газа~\eqref{eq:free_gas_rate}. Его правая часть
равна~$\mdot/\rho_g^{\text{ст}}$ (источник свободного газа за~счёт
дегазации). Следовательно:
\begin{multline}\label{eq:hg_components_derived}
  \pdt{}\!\left(\phi\,\rho_g^{\text{ст}}\,
    \frac{S_g}{B_g}\,h_g\right)
  + \divr{\rho_g^{\text{ст}}\,
    \frac{u_g}{B_g}\,h_g} =
    \\
    +\mdot\,h_g
  +
    \left(\phi\,\rho_g^{\text{ст}}\,
      \frac{S_g}{B_g}\right)
    \pdt{h_g}\!
    + \left(\rho_g^{\text{ст}}\,
      \frac{u_g}{B_g}\right)\pdr{h_g}.
\end{multline}

Сумма уравнений~\eqref{eq:F_components_derived} и~\eqref{eq:hg_components_derived}
содержит вклад фазового перехода
\begin{equation}\label{eq:phase_change_heat}
  (h_g - \bar{h}_{go})\,\mdot \equiv L\,\mdot,
\end{equation}
где введена \textbf{удельная теплота дегазации}:
\begin{equation}\label{eq:L_def}
  \boxed{
    L(p,T) = h_g(p,T) - \bar{h}_{go}(p, T, R_s^{\text{eq}}(p,T)).
  }
\end{equation}

Парциальная энтальпия~$\bar{h}_{go}$ определена для произвольного
состава нефтяной фазы: $\bar{h}_{go} = \bar{h}_{go}(p, T, R_s)$,
где~$R_s$~--- фактическое содержание газа в~нефти (не~обязательно
равновесное). В~определении~$L$ состав фиксируется при
\textbf{равновесном} значении $R_s = R_s^{\text{eq}}(p,T)$,
обеспечивающем сосуществование нефтяной и~газовой фаз.
Именно при этом составе имеет смысл говорить о~переходе молекул
газа между фазами. Поэтому $L$ выражается как функция только~$p$
и~$T$.

В случае $L > 0$ свободный газ имеет большую энтальпию, чем растворённый.
При выделении газа ($\mdot > 0$) и~$L > 0$: система охлаждается~---
энергия расходуется на~перевод молекул газа из~растворённого состояния
в~свободное.

\paragraph{Поведение в~однофазной зоне.}
В~однофазной зоне ($S_g = 0$, режим~2) газосодержание~$R_s$
определяется из баланса масс. При отсутствии дополнительных
источников газа и~постоянном составе притока $R_s = R_s^0 =
\mathrm{const}$, поэтому $\partial R_s/\partial t = 0$
и~$\partial R_s/\partial r = 0$. Из~определения~\eqref{eq:mdot}
следует $\mdot = 0$, и~вклад дегазации $Q_{\mathrm{deg}}
= -L\,\mdot = 0$~--- дегазации нет.

Формально в~определении~\eqref{eq:L_def} фигурирует энтальпия
газовой фазы~$h_g$, которая не~имеет физического смысла при~$S_g=0$
(газовой фазы не~существует). Однако при~$\mdot = 0$ произведение
$L\,\mdot = 0$ независимо от~значения~$L$, поэтому значение~$L$
в~однофазной зоне несущественно и~не~влияет на~решение.

\paragraph{Вклады изменения энтальпий при фиксированном составе.}

В~уравнениях~\eqref{eq:F_components_derived}
и~\eqref{eq:hg_components_derived} помимо членов фазового перехода
($-\mdot\,\bar{h}_{go}$ и~$+\mdot\,h_g$) остались слагаемые
с~производными энтальпий по~$t$ и~$r$. Перечислим их явно.

По~аналогии с~\eqref{eq:F_components_der} для нефтяного компонента $\bar{h}_{oo}$
(первое слагаемое в~\eqref{eq:H_fluids}, \eqref{eq:F_components}):
\begin{multline}\label{eq:hoo_terms}
  \pdt{}\!\left(\phi\,\rho_o^{\text{ст}}\,
    \frac{S_o}{B_o}\,\bar{h}_{oo}\right)
  + \divr{\rho_o^{\text{ст}}\,
    \frac{u_o}{B_o}\,\bar{h}_{oo}} =
    \\
    \bar{h}_{oo}\!\left[
    \pdt{}\!\left(\phi\,\rho_o^{\text{ст}}\,
      \frac{S_o}{B_o}\right)
    + \divr{\rho_o^{\text{ст}}\,\frac{u_o}{B_o}}
  \right]
  +
  \\
    \left(\phi\,\rho_o^{\text{ст}}\,
      \frac{S_o}{B_o}\right)
    \pdt{\bar{h}_{oo}}\!
    + \left(\rho_o^{\text{ст}}\,
      \frac{u_o}{B_o}\right)\pdr{\bar{h}_{oo}}.
\end{multline}
Выражение в~квадратных скобках представляет собой плотность $\rho_o^{\text{ст}}$,
умноженную на~левую
часть баланса нефтяного компонента~\eqref{eq:mass_oil}, которая
равна нулю. Следовательно, первый член обращается в~нуль, и~остаются
только производные~$\bar{h}_{oo}$
\begin{equation}\label{eq:hoo_remainder}
    \pdt{}\!\left(\phi\,\rho_o^{\text{ст}}\,
    \frac{S_o}{B_o}\,\bar{h}_{oo}\right)
  + \divr{\rho_o^{\text{ст}}\,
    \frac{u_o}{B_o}\,\bar{h}_{oo}} =\left(\phi\,\rho_o^{\text{ст}}\,
      \frac{S_o}{B_o}\right)
    \pdt{\bar{h}_{oo}}\!
    + \left(\rho_o^{\text{ст}}\,
      \frac{u_o}{B_o}\right)\pdr{\bar{h}_{oo}}.
\end{equation}

Таким образом, уравнение~\eqref{eq:energy_cons} принимает следующий вид
с~учётом выполненных преобразований
\begin{multline}\label{eq:energy_cons_der}
  \pdt{}\left( (1 - \phi)\,(\rho c_p)_s\,T - \phi\,p\right)
      +
      \\
    \left(\phi\,\rho_g^{\text{ст}}\,
      \frac{S_g}{B_g}\right)
    \pdt{h_g}\!
    + \left(\rho_g^{\text{ст}}\,
      \frac{u_g}{B_g}\right)\pdr{h_g}  +
      \\
    \left(\phi\,\rho_g^{\text{ст}}\,
      \frac{R_s\,S_o}{B_o}\right)
    \pdt{\bar{h}_{go}}\!
    + \left(\rho_g^{\text{ст}}\,
      \frac{R_s\,u_o}{B_o}\right)\pdr{\bar{h}_{go} } +
      \\
  \left(\phi\,\rho_o^{\text{ст}}\,
      \frac{S_o}{B_o}\right)
    \pdt{\bar{h}_{oo}}\!
    + \left(\rho_o^{\text{ст}}\,
      \frac{u_o}{B_o}\right)\pdr{\bar{h}_{oo}} =
  \\
  \divr{\leff\,\pdr{T}} - L\mdot,
\end{multline}


Уравнение~\eqref{eq:energy_cons_der} содержит производные парциальных
энтальпий~$\bar{h}_{oo}$, $\bar{h}_{go}$ и~энтальпии газовой
фазы~$h_g$. Парциальные энтальпии~$\bar{h}_{oo}$ и~$\bar{h}_{go}$ зависят
от~$p$, $T$ и~состава~$R_s$. Их полные производные по~$t$
содержат члены с~$\partial R_s/\partial t$:
\[
  \pdt{\bar{h}_{oo}} = \left.\pdt{\bar{h}_{oo}}\right|_{R_s}
    + \frac{\partial\bar{h}_{oo}}{\partial R_s}\,\pdt{R_s},
  \qquad
  \pdt{\bar{h}_{go}} = \left.\pdt{\bar{h}_{go}}\right|_{R_s}
    + \frac{\partial\bar{h}_{go}}{\partial R_s}\,\pdt{R_s}.
\]
Запишем соотношение Гиббса--Дюгема для бинарной смеси (нефтяная фаза)
при постоянных~$p$ и~$T$
\begin{equation}\label{eq:gibbs_duhem}
  \rho_o^{\text{ст}}\,d\bar{h}_{oo}\big|_{p,T}
  + R_s\,\rho_g^{\text{ст}}\,d\bar{h}_{go}\big|_{p,T} = 0.
\end{equation}
В~результате группировки слагаемых в~уравнении~\eqref{eq:energy_cons_der}
с~учётом соотношения~\eqref{eq:gibbs_duhem} запишем
\begin{multline}\label{eq:energy_cons_derived}
  \pdt{}\left( (1 - \phi)\,(\rho c_p)_s\,T - \phi\,p\right)
      +
      \\
    \left(\phi\,\rho_g^{\text{ст}}\,
      \frac{S_g}{B_g}\right)
    \pdt{h_g}\!
    + \left(\rho_g^{\text{ст}}\,
      \frac{u_g}{B_g}\right)\pdr{h_g}  +
      \\
    \phi\,\frac{S_o}{B_o}
    \left.\pdt{}\left(\rho_g^{\text{ст}}\,R_s\,\bar{h}_{go} + \rho_o^{\text{ст}}\,\bar{h}_{oo}\right)\right|_{R_s}
    +
      \frac{u_o}{B_o}\left.\pdr{}\left(\rho_g^{\text{ст}}\,R_s\,\bar{h}_{go} + \rho_o^{\text{ст}}\,\bar{h}_{oo}\right) \right|_{R_s} =
  \\
  \divr{\leff\,\pdr{T}} - L\mdot,
\end{multline}
Преобразуем уравнение~\eqref{eq:energy_cons_derived} далее с~учётом
определения~\eqref{eq:rho_oil} плотности нефтяной фазы и
определения энтальпии нефти~\eqref{eq:oil_enthalpy_components}
\begin{multline}\label{eq:energy_cons_derived2}
  \pdt{}\left( (1 - \phi)\,(\rho c_p)_s\,T - \phi\,p\right)
      +
      \\
    \phi\,\rho_g^{\text{ст}}\,
      \frac{S_g}{B_g}
    \pdt{h_g}\!
    + \rho_g^{\text{ст}}\,
      \frac{u_g}{B_g}\pdr{h_g}  +
      \\
    \phi\,\frac{S_o}{B_o}\,(\rho_o^{\text{ст}}+\rho_g^{\text{ст}}\,R_s)
    \left.\pdt{h_o}\right|_{R_s}
    +
      \frac{u_o}{B_o}\,(\rho_o^{\text{ст}}+\rho_g^{\text{ст}}\,R_s)\left.\pdr{h_o}\right|_{R_s} =
  \\
  \divr{\leff\,\pdr{T}} - L\mdot,
\end{multline}

Дифференциал энтальпии нефтяной фазы при фиксированном
составе принимает следующий вид
\begin{equation}\label{eq:dh_oil}
  dh_o\big|_{R_s} = c_{p,o}\,dT + v_o\,(1 - \alpha_o\,T)\,dp,
\end{equation}
где $v_o = 1/\rho_o$~--- удельный объём,
$c_{p,o}$~--- удельная изобарная теплоёмкость (двухкомпонентной) нефти,
$\alpha_o$~--- её коэффициент объёмного теплового расширения.
Первое слагаемое характеризует изменение энтальпии при изобарном нагреве,
второе~--- при изотермическом изменении давления. Аналогичное соотношение записывается
для газовой фазы.

Подставляя~\eqref{eq:dh_oil} в~уравнение~\eqref{eq:energy_cons_derived2}, получим
\begin{multline}\label{eq:energy_cons_derived3}
  \pdt{}\left( (1 - \phi)\,(\rho c_p)_s\,T - \phi\,p\right)
      +
      \\
    \phi\,\rho_g\,S_g
    \left(c_{p,g}\,\pdt{T} + v_g\,(1 - \alpha_g\,T)\,\pdt{p}\right)\!
    + \rho_g\,u_g\left(c_{p,g}\,\pdr{T} + v_g\,(1 - \alpha_g\,T)\,\pdr{p}\right)  +
    \\
    \phi\,S_o\,\rho_o
    \left(c_{p,o}\,\pdt{T} + v_o\,(1 - \alpha_o\,T)\,\pdt{p}\right)
    +
      u_o\,\rho_o\left(c_{p,o}\,\pdr{T} + v_o\,(1 - \alpha_o\,T)\,\pdr{p}\right) =
  \\
  \divr{\leff\,\pdr{T}} - L\mdot.
\end{multline}

Раскроем производные от слагаемых, описывающих теплоёмкость пород и
работу внешних сил
\begin{equation}\label{eq:rock_dt}
  \pdt{}\!\left[
    (1-\phi)\,(\rho c_p)_s\,T-\phi\,p
  \right]
  = (1-\phi)\,(\rho c_p)_s\,\pdt{T}
    - (\rho c_p)_s\,T\,\phi\,c_r\,\pdt{p}-\phi\,(1 + p\,c_r)\,\pdt{p}.
\end{equation}
Подставим это выражение в~уравнение~\eqref{eq:energy_cons_derived3}
и~перегруппируем слагаемые
\begin{multline}\label{eq:energy_cons_derived4}
  \left[(1-\phi)\,(\rho c_p)_s\, +\phi\,\rho_g\,S_gc_{p,g}\,  +\phi\,S_o\,\rho_oc_{p,o}\right]\pdt{T}
    +
    \\
    (\phi\,S_g\,(1 - \alpha_g\,T)+\phi\,S_o\,(1 - \alpha_o\,T)\,-(\rho c_p)_s\,T\,\phi\,c_r-\phi\,(1 + p\,c_r))\,\pdt{p}
      +
      \\
    (\rho_g\,u_gc_{p,g}+\rho_ou_oc_{p,o})\,\pdr{T} + (u_g\,(1 - \alpha_g\,T)+ u_o\,(1 - \alpha_o\,T))\pdr{p} =
  \\
  \divr{\leff\,\pdr{T}} - L\mdot.
\end{multline}

Определим эффективные характеристики гомогенной смеси
\begin{equation}\label{eq:1}
\rhocpeff
    = (1-\phi)\,(\rho c_p)_s
      + \phi\,(S_o\,\rho_o\,c_{p,o} + S_g\,\rho_g\,c_{p,g}),
\end{equation}
\begin{equation}\label{eq:2}
\rhocpconv
    = \rho_o\,c_{p,o}\,u_o + \rho_g\,c_{p,g}\,u_g,
\end{equation}
\begin{equation}\label{eq:3}
\Qad = \phi\left[T\,\sum_\alpha S_\alpha\alpha_\alpha + (\rho c_p)_s\,T\,\,c_r + p\,c_r\right]\pdt{p},
\end{equation}
\begin{equation}\label{eq:4}
\QJT = -\left[u_g\,(1 - \alpha_g\,T)+ u_o\,(1 - \alpha_o\,T)\right]\pdr{p},
\end{equation}
\begin{equation}\label{eq:5}
\Qdeg = - L(p,T)\mdot
\end{equation}
и~получим уравнение баланса энергии
\begin{equation}\label{eq:energy_final}
  \boxed{
    \rhocpeff\,\pdt{T}
    + \rhocpconv\,\pdr{T}
    = \divr{\leff\,\pdr{T}}
    + \QJT + \Qad + \Qdeg.
  }
\end{equation}

\paragraph{Запись через коэффициент Джоуля--Томсона.}

Введём \textbf{коэффициент Джоуля--Томсона}~--- изменение
температуры при изоэнтальпийном расширении:
\begin{equation}\label{eq:muJT_def}
  \mu_{\mathrm{JT},\alpha}
  = \left.\frac{\partial T}{\partial p}\right|_{h_\alpha}
  = \frac{v_\alpha\,(\alpha_\alpha\,T - 1)}{c_{p,\alpha}},
  \qquad \alpha \in \{o,\,g\},
\end{equation}
где $v_\alpha = 1/\rho_\alpha$. Отсюда
\begin{equation}\label{eq:alpha_via_muJT}
  1 - \alpha_\alpha\,T
  = -\rho_\alpha\,c_{p,\alpha}\,\mu_{\mathrm{JT},\alpha}.
\end{equation}

Подставим~\eqref{eq:alpha_via_muJT} в~определения
$\QJT$~\eqref{eq:4} и~$\Qad$~\eqref{eq:3}
\begin{equation}\label{eq:QJT_muJT}
  \boxed{
    \QJT^*
    = \sum_\alpha \rho_\alpha\,c_{p,\alpha}\,
      \mu_{\mathrm{JT},\alpha}\,u_\alpha\,\pdr{p}.
  }
\end{equation}
\begin{equation}\label{eq:Qad_muJT}
  \boxed{
    \Qad^*
    = \phi\!\left[
      \sum_\alpha S_\alpha\!\left(1
        + \rho_\alpha\,c_{p,\alpha}\,\mu_{\mathrm{JT},\alpha}\right)
      + c_r\!\left(p + (\rho c_p)_s\,T\right)
    \right]\pdt{p}.
  }
\end{equation}
Здесь использовано $\alpha_\alpha\,T
= 1 + \rho_\alpha\,c_{p,\alpha}\,\mu_{\mathrm{JT},\alpha}$
из~\eqref{eq:alpha_via_muJT}.

В~PVT-контракте (раздел~\ref{ssec:closures}) коэффициент
$\mu_{\mathrm{JT},\alpha}(p,T)$ заменяет~$\alpha_\alpha(p,T)$
как входной параметр. Пересчёт, при необходимости:
$\alpha_\alpha = (1 + \rho_\alpha\,c_{p,\alpha}\,
\mu_{\mathrm{JT},\alpha})/T$.


%=============================================================================
\section{Раскрытие балансов масс и уравнение давления}\label{app:expansion}
%=============================================================================

Ниже приведены промежуточные выкладки, полезные для анализа
структуры уравнений. Эти формулы \textbf{не используются} при построении
численной схемы~--- дискретизация основана на~исходных консервативных
уравнениях~\eqref{eq:mass_oil}, \eqref{eq:mass_gas}.

Сжимаемость порового пространства
$c_r = (1/\phi)\,d\phi/dp$; следствие:
\begin{equation}\label{eq:dphi_dt}
  \pdt{\phi} = \phi\,c_r\,\pdt{p},
  \qquad
  \pdr{\phi} = \phi\,c_r\,\pdr{p}.
\end{equation}

\subsection{Раскрытие баланса нефтяного компонента}

Раскроем временную производную в~\eqref{eq:mass_oil}:
\begin{equation}\label{eq:oil_dt_step1}
  \pdt{}\!\left(\phi\,\xio\right)
  = \phi\,\pdt{\xio}
    + \xio\,\pdt{\phi}.
\end{equation}

С~учётом~\eqref{eq:dphi_dt}:
$\xio\,\partial\phi/\partial t = \phi\,c_r\,\xio\,\partial p/\partial t$.

Из~$\xio = (1 - S_g)/B_o$:
\begin{equation}\label{eq:dxi_dt}
  \pdt{\xio}
  = -\frac{1}{B_o}\,\pdt{S_g}
    - \frac{\xio}{B_o}\,\pdt{B_o}.
\end{equation}

Полная производная $B_o(p,T,R_s)$:
\begin{equation}\label{eq:dBo_dt}
  \frac{1}{B_o}\,\pdt{B_o}
  = -c_o\,\pdt{p}
    + \beta_o\,\pdt{T}
    + \frac{1}{B_o}\left.\frac{\partial B_o}{\partial R_s}\right|_{p,T}
      \!\left(
        \frac{\partial R_s}{\partial p}\,\pdt{p}
        + \frac{\partial R_s}{\partial T}\,\pdt{T}
      \right).
\end{equation}

Введём <<полные>> коэффициенты вдоль кривой насыщения:
\begin{align}
  \csat &= c_o
    - \frac{1}{B_o}\left.\frac{\partial B_o}{\partial R_s}\right|_{p,T}
      \left.\frac{\partial R_s}{\partial p}\right|_T,
      \label{eq:co_sat} \\[4pt]
  \bsat &= \beta_o
    + \frac{1}{B_o}\left.\frac{\partial B_o}{\partial R_s}\right|_{p,T}
      \left.\frac{\partial R_s}{\partial T}\right|_p.
      \label{eq:beta_o_sat}
\end{align}

$\csat < c_o$: при увеличении давления нефть поглощает газ и~набухает.

Коэффициенты $c_o$, $\beta_o$, $\alpha_o$~--- при \textbf{фиксированном}
составе (для уравнения энергии); $\csat$, $\bsat$~--- \textbf{вдоль
кривой насыщения} (для балансов масс).

Раскрытая форма баланса нефтяного компонента:
\begin{equation}\label{eq:mass_oil_expanded}
  -\frac{\phi}{B_o}\,\pdt{S_g}
  + \phi\,\xio\,(c_r + \csat)\,\pdt{p}
  - \phi\,\xio\,\bsat\,\pdt{T}
  + \divr{\frac{u_o}{B_o}} = 0.
\end{equation}


\subsection{Раскрытие баланса газового компонента}

Обозначим $\xi_g = R_s\,\xio + S_g/B_g$. Раскрытая форма:
\begin{multline}\label{eq:mass_gas_expanded}
  \phi\!\left(\frac{1}{B_g} - \frac{R_s}{B_o}\right)
    \pdt{S_g}
  + \phi\!\left[\xio\,\frac{\partial R_s}{\partial p}
    + \xio\,R_s\,(c_r + \csat)
    + \frac{S_g}{B_g}\,(c_g + c_r)\right]
    \pdt{p}
  \\
  + \phi\!\left[\xio\,\frac{\partial R_s}{\partial T}
    - R_s\,\xio\,\bsat
    - \frac{S_g}{B_g}\,\beta_g\right]
    \pdt{T}
  + \divr{\left[
    \frac{R_s\,u_o}{B_o} + \frac{u_g}{B_g}
  \right]} = 0.
\end{multline}


\subsection{Уравнение давления}

Исключая $\partial S_g/\partial t$ из~\eqref{eq:mass_oil_expanded}
и~\eqref{eq:mass_gas_expanded}, получаем:
\begin{equation}\label{eq:pressure}
  C_p\,\pdt{p}
  + C_T\,\pdt{T}
  + \mathcal{D} = 0,
\end{equation}
где
\begin{equation}\label{eq:Cp}
  C_p = \phi\!\left[
    S_o\,(c_r + \csat)\!\left(\frac{1}{B_g} - \frac{R_s}{B_o}\right)
    + \xio\!\left(\frac{\partial R_s}{\partial p}
      + R_s\,(c_r + \csat)\right)
    + \frac{S_g}{B_g}\,(c_g + c_r)
  \right],
\end{equation}
\begin{equation}\label{eq:CT}
  C_T = \phi\!\left[
    -S_o\,\bsat\!\left(\frac{1}{B_g} - \frac{R_s}{B_o}\right)
    + \xio\!\left(\frac{\partial R_s}{\partial T}
      - R_s\,\bsat\right)
    - \frac{S_g}{B_g}\,\beta_g
  \right],
\end{equation}
\begin{equation}\label{eq:D_operator}
  \mathcal{D} = \frac{B_o}{B_g}\,\divr{\frac{u_o}{B_o}}
    + \frac{u_o}{B_o}\,\pdr{R_s}
    + \divr{\frac{u_g}{B_g}}.
\end{equation}

Множитель $B_o/B_g$ перед первой дивергенцией~--- следствие исключения
$\partial S_g/\partial t$ из~двух балансов масс.

Эволюция газонасыщенности (из~\eqref{eq:mass_oil_expanded}):
\begin{equation}\label{eq:Sg_ode}
  \pdt{S_g}
  = S_o\,(c_r + \csat)\,\pdt{p}
    - S_o\,\bsat\,\pdt{T}
    + \frac{B_o}{\phi}\,\divr{\frac{u_o}{B_o}}.
\end{equation}


\end{document}
